\documentclass[../main.tex]{subfiles}
\graphicspath{{\subfix{../image/}}{\subfix{../figure/}}}

\begin{document}

\chapter{拟线性抛物方程}\label{chapter:ch06}

本章研究初边值问题
\begin{align*}
\begin{cases}
\partial_t u = \Delta\phi(u) & \text{在 }\Omega\times(0,T)=:Q,\\
u=0 & \text{在 }\partial\Omega\times(0,T)=:\Gamma,\\
u=u_0 & \text{在 }\Omega\times\{0\},
\end{cases}
\qquad(P)
\end{align*}
其中 $\Omega\subseteq\mathbb{R}^N$（$N\geqslant3$）为具有光滑边界 $\partial\Omega$ 的有界开集，$T>0$，$u_0$ 为有限 Radon 测度，且
\begin{align*}
u_0\in\mathscr{R}_f(\Omega). \qquad(A0)
\end{align*}
关于扩散函数 $\phi:\mathbb{R}\to\mathbb{R}$，我们采用如下假设：
\begin{align*}
\phi\in C_b(\mathbb{R})\cap C^\infty(\mathbb{R}) \text{ 在 }\mathbb{R}\text{ 上非减且非常数，}\phi(0)=0; \qquad(A1)
\end{align*}
\begin{align*}
\phi\in W^{1,\infty}(\mathbb{R}). \qquad(A2)
\end{align*}
更一般的假设将在\cref{section:ch06-无界phi}中使用。

由假设 (A1)，存在
\begin{align}
\lim_{s\to\pm\infty}\phi(s)=:\phi_{\pm\infty}\in\mathbb{R}. \label{align:ch6-6.1}
\end{align}
若 $\phi'$ 在 $\pm\infty$ 处有极限，则该极限为零，因而问题 (P) 在 $\pm\infty$ 处是退化抛物的。

此后对任意 $u\in\mathscr{R}_f(Q)$，我们记 $u_r\in L^1(Q)$ 与 $u_s\in\mathscr{R}_f(Q)$ 分别为 $u$ 关于 Lebesgue 测度的绝对连续部分 $u_{\mathrm{ac}}$ 与奇异部分的密度（见\cref{definition:ch01-Lebesgue分解}）。空间 $\mathscr{R}_f(\Omega)$ 中使用类似记号，例如对 $u_0$。对任意 $u\in\mathscr{R}_f(Q)$，我们用 $u(\cdot,t)$（$t\in(0,T)$）记其离散分解 $\{u_t\}_{t\in(0,T)}$（见\cref{definition:ch05-离散分解定义}）。"几乎处处"始终指关于 Lebesgue 测度。为简便记 $dt\equiv\mathrm{d}\lambda(t)$，$dx\equiv\mathrm{d}\lambda_N(x)$，$dx\,\mathrm{d}t\equiv\mathrm{d}\lambda_{N+1}(x,t)$。如通常，$\langle\cdot,\cdot\rangle$ 表示 $\mathscr{R}_f(\Omega)$ 与 $C_0(\Omega)$ 之间的对偶映射。按记号滥用，对 $\rho\in C_b(\overline\Omega)$ 也置
\begin{align*}
\langle\mu,\rho\rangle\equiv\int_\Omega\rho\,\mathrm{d}\mu.
\end{align*}

\section{问题陈述与预备结果}\label{section:ch06-问题陈述}

\subsection{弱解}\label{subsection:ch06-弱解}

下面引入我们的主要解概念。

\begin{definition}\label{definition:ch06-弱解}
设 (A0)--(A1) 成立。问题 (P) 的弱解指任意 $u\in L^\infty_{w^*}(0,T;\mathscr{R}_f(\Omega))$，满足：
\begin{enumerate}[(i)]
\item $\phi(u_r)\in L^2(0,T;H^1_0(\Omega))$;
\item 对任意 $\zeta\in C^1([0,T];C^1_c(\Omega))$ 且 $\zeta(\cdot,T)=0$ 在 $\Omega$ 上，有
\begin{align}
\int_0^T\langle u(\cdot,t),\partial_t\zeta(\cdot,t)\rangle\,\mathrm{d}t
= \iint_Q\nabla\phi(u_r)\cdot\nabla\zeta\,\mathrm{d}x\,\mathrm{d}t
-\langle u_0,\zeta(\cdot,0)\rangle. \label{align:ch6-6.2}
\end{align}
\end{enumerate}
\end{definition}

置 $C^{1,2}(Q):=\{\zeta\in C(Q)\mid\exists\,\partial_t\zeta,\partial_{x_i}\zeta,\partial^2_{x_ix_j}\zeta\in C(Q)\ \forall i,j=1,\dots,N\}$。更弱的解概念如下。

\begin{definition}\label{definition:ch06-极弱解}
设 (A0)--(A1) 成立。(P) 的极弱解指任意 $u\in L^\infty_{w^*}(0,T;\mathscr{R}_f(\Omega))$，满足
\begin{align}
\int_0^T\langle u(\cdot,t),\partial_t\zeta(\cdot,t)\rangle\,\mathrm{d}t
= -\iint_Q\phi(u_r)\Delta\zeta\,\mathrm{d}x\,\mathrm{d}t
-\langle u_0,\zeta(\cdot,0)\rangle \label{align:ch6-6.3}
\end{align}
其中 $\zeta\in C^{1,2}(Q)$ 在 $\Gamma\cup(\Omega\times\{T\})$ 上为零。
\end{definition}

\begin{remark}\label{remark:ch06-6.1.1}
由标准逼近论证，可在\eqref{align:ch6-6.2}中取在 $\Gamma\cup(\Omega\times\{T\})$ 上为零的试验函数 $\zeta\in C^1(Q)$。因此 (P) 的每个弱解都是极弱解。
\end{remark}

如下命题的 (ii) 款说明 (P) 中初值条件如何满足。

\begin{proposition}\label{proposition:ch06-6.1.1}\label{proposition:ch6-1.1}
设 (A0)--(A1) 成立，$u$ 为问题 (P) 的极弱解。则：
\begin{enumerate}[(i)]
\item 存在零测集 $F^*\subseteq(0,T)$，使对每个 $t\in(0,T)\setminus F^*$ 及所有 $\rho\in C^2_c(\Omega)$，
\begin{align}
\langle u(\cdot,t),\rho\rangle-\langle u_0,\rho\rangle
= \int_0^t\mathrm{d}s\int_\Omega\phi(u_r(x,s))\Delta\rho(x)\,\mathrm{d}x; \label{align:ch6-6.4}
\end{align}
\item 对任意 $\rho\in C_0(\Omega)$，有
\begin{align}
\operatorname*{ess\,lim}_{t\to0^+}\langle u(\cdot,t),\rho\rangle=\langle u_0,\rho\rangle, \label{align:ch6-6.5a}\\
\operatorname*{ess\,lim}_{t\to t_0}\langle u(\cdot,t),\rho\rangle=\langle u(\cdot,t_0),\rho\rangle
\quad\text{对几乎所有 }t_0\in(0,T). \label{align:ch6-6.5b}
\end{align}
\end{enumerate}
\end{proposition}

\begin{proof}
(i) 我们将证明存在零测集 $F^*\subseteq(0,T)$，使对所有 $\rho\in C_b(\overline\Omega)$ 及 $t\in(0,T)\setminus F^*$，
\begin{align}
\lim_{q\to\infty}\left(2q\int_{t-\frac{1}{q}}^{t+\frac{1}{q}}
\left|\langle u(\cdot,s),\rho\rangle-\langle u(\cdot,t),\rho\rangle\right|\mathrm{d}s\right)=0. \label{align:ch6-6.6}
\end{align}
利用\eqref{align:ch6-6.3}与\eqref{align:ch6-6.6}，可证等式\eqref{align:ch6-6.4}。设 $\rho\in C^2_c(\Omega)$，$t_1\in(0,T)\setminus F^*$。由标准正则化论证，可在\eqref{align:ch6-6.3}中取 $\zeta(x,t)=\rho(x)k_q(t)$，其中 $q\in\mathbb{N}$，$q\geqslant1/(T-t_1)+1$，且
\begin{align}
k_q(t):=\chi_{(0,t_1]}(t)+[q(t_1-t)+1]\chi_{(t_1,t_1+\frac{1}{q}]}(t) \quad(t\in[0,T]). \label{align:ch6-6.7}
\end{align}
由\eqref{align:ch6-6.3}得
\begin{align*}
q\int_{t_1}^{t_1+\frac{1}{q}}\langle u(\cdot,t),\rho\rangle\,\mathrm{d}t-\langle u_0,\rho\rangle
= \iint_Q\phi(u_r)\Delta\rho(x)k_q(t)\,\mathrm{d}x\,\mathrm{d}t.
\end{align*}
由于对所有 $t\in(0,T)$ 有 $\lim_{q\to\infty}k_q(t)=\chi_{(0,t_1]}(t)$，令 $q\to\infty$ 并利用\eqref{align:ch6-6.6}，得
\begin{align*}
\langle u(\cdot,t_1),\rho\rangle-\langle u_0,\rho\rangle
= \int_0^{t_1}\mathrm{d}t\int_\Omega\phi(u_r(x,t))\Delta\rho(x)\,\mathrm{d}x,
\end{align*}
即\eqref{align:ch6-6.4}。

为证\eqref{align:ch6-6.6}，注意因 $u\in L^\infty_{w^*}(0,T;\mathscr{R}_f(\Omega))$，由\cref{theorem:ch05-弱星可测}存在零测集 $F_0\subseteq(0,T)$，使对每个 $t\in(0,T)\setminus F_0$，离散分解 $u(\cdot,t)\in\mathscr{R}_f(\Omega)$ 有定义；且对任意 $\rho\in C_b(\overline\Omega)$，映射 $t\mapsto\psi_\rho(t):=\langle u(\cdot,t),\rho\rangle$ 属于 $L^\infty(0,T)$。设 $D\equiv\{\rho_k\}$ 为 $C_b(\overline\Omega)$ 的可数稠密集。因 $\psi_{\rho_k}\in L^\infty(0,T)$，存在零测集 $F_k\subseteq(0,T)$，使对每个 $t\in(0,T)\setminus F_k$，\eqref{align:ch6-6.6}成立（见\cref{corollary:ch02-平均收敛}(ii)）。不妨设 $F_0\subseteq F_k$ 对每个 $k\in\mathbb{N}$。置 $F^*:=\bigcup_{k=1}^\infty F_k$。则对所有 $t\in(0,T)\setminus F^*$ 及 $k\in\mathbb{N}$，有
\begin{align*}
\lim_{q\to\infty}\left(2q\int_{t-\frac{1}{q}}^{t+\frac{1}{q}}
\left|\langle u(\cdot,s),\rho_k\rangle-\langle u(\cdot,t),\rho_k\rangle\right|\mathrm{d}s\right)=0.
\end{align*}
由 $D$ 在 $C_b(\overline\Omega)$ 中的稠密性，等式\eqref{align:ch6-6.6}随之成立。这证明了该断言。

(ii) 设 $\{t_n\}\subseteq(0,T)\setminus F^*$，$t_n\to0^+$（$n\to\infty$）。因 $\phi(u_r)\in L^\infty(Q)$，在\eqref{align:ch6-6.4}中取 $t=t_n$ 并令 $n\to\infty$，得
\begin{align*}
\langle u(\cdot,t_n),\rho\rangle\to\langle u_0,\rho\rangle
\quad\text{对所有 }\rho\in C^2_c(\Omega).
\end{align*}
又因 $u\in L^\infty_{w^*}(0,T;\mathscr{R}_f(\Omega))$，有 $\sup_{n\in\mathbb{N}}\|u(\cdot,t_n)\|_{\mathscr{R}_f(\Omega)}<\infty$。由 Banach--Alaoglu 定理，存在 $\mu_0\in\mathscr{R}_f(\Omega)$ 及子列 $\{t_{n_k}\}\subseteq\{t_n\}$，使 $u(\cdot,t_{n_k})\stackrel{*}{\rightharpoonup}\mu_0$ 在 $\mathscr{R}_f(\Omega)$。由标准稠密性论证得 $\mu_0=u_0$ 在 $\mathscr{R}_f(\Omega)$。故 $u(\cdot,t_n)\stackrel{*}{\rightharpoonup}u_0$ 沿整列 $\{t_n\}$，\eqref{align:ch6-6.5a}由 $\{t_n\}$ 的任意性得证。对等式
\begin{align*}
\langle u(\cdot,t),\rho\rangle-\langle u(\cdot,t_0),\rho\rangle
= \int_{t_0}^t\mathrm{d}s\int_\Omega\phi(u_r(x,s))\Delta\rho(x)\,\mathrm{d}x
\end{align*}
（这是\eqref{align:ch6-6.4}对几乎所有 $t_0,t\in(0,T)$ 的直接推论）应用同样论证，即证\eqref{align:ch6-6.5b}。结论随之成立。
\end{proof}

\begin{remark}\label{remark:ch06-6.1.2}
由\cref{proposition:ch06-6.1.1}易见，映射 $t\mapsto u(\cdot,t)$ 存在一个对所有 $t\in[0,T]$ 有定义的代表元，使 $u(\cdot,0)=u_0$，且对每个 $\rho\in C_0(\Omega)$，
\begin{align}
\lim_{t\to0^+}\langle u(\cdot,t),\rho\rangle=\langle u_0,\rho\rangle, \label{align:ch6-6.8a}\\
\lim_{t\to t_0}\langle u(\cdot,t),\rho\rangle=\langle u(\cdot,t_0),\rho\rangle
\quad\text{对所有 }t_0\in[0,T]. \label{align:ch6-6.8b}
\end{align}
此后凡陈述映射 $t\mapsto u(\cdot,t)$ 在 $[0,T]$ 某子区间上对所有 $t$ 成立的性质时，均指此连续代表元。此时等式\eqref{align:ch6-6.4}对每个 $t\in(0,T)$ 成立，等式\eqref{align:ch6-6.5a}由\eqref{align:ch6-6.8a}替代。
\end{remark}

\subsection{弱熵解}\label{subsection:ch06-弱熵解}

为后文需要，引入问题 (P) 的另一类解。设 (A1) 成立。对每个非减 $f\in C(\mathbb{R})$，置
\begin{align}
F_\phi(y):=\int_0^y f(\phi(z))\,\mathrm{d}z \quad(y\in\mathbb{R}), \label{align:ch6-6.9}
\end{align}
\begin{align}
M^\pm_{F_\phi}:=\lim_{y\to\pm\infty}\frac{F_\phi(y)}{y}
=\lim_{y\to\pm\infty}f(\phi(y))\in\mathbb{R}. \label{align:ch6-6.10}
\end{align}

\begin{definition}\label{definition:ch06-熵解}
设 (A0)--(A1) 成立。问题 (P) 的熵解指 (P) 的任意弱解 $u$，满足熵不等式：
\begin{align}
\begin{aligned}
&\iint_Q\{F_\phi(u_r)\partial_t\zeta-f(\phi(u_r))\nabla\phi(u_r)\cdot\nabla\zeta\}\,\mathrm{d}x\,\mathrm{d}t\\
&\quad+M^+_{F_\phi}\int_0^T\langle u^+_s(\cdot,t),\partial_t\zeta(\cdot,t)\rangle\,\mathrm{d}t
-M^-_{F_\phi}\int_0^T\langle u^-_s(\cdot,t),\partial_t\zeta(\cdot,t)\rangle\,\mathrm{d}t\\
&\geqslant -\int_\Omega F_\phi(u_{0r})\zeta(x,0)\,\mathrm{d}x
-M^+_{F_\phi}\langle u^+_{0s},\zeta(\cdot,0)\rangle
+M^-_{F_\phi}\langle u^-_{0s},\zeta(\cdot,0)\rangle
\end{aligned} \label{align:ch6-6.11}
\end{align}
对任意非减 $f\in C(\mathbb{R})$ 及 $\zeta\in C^1([0,T];C^1_c(\Omega))$，$\zeta\geqslant0$，$\zeta(\cdot,T)=0$ 在 $\Omega$ 上成立。
\end{definition}

易见\eqref{align:ch6-6.11}中各项均有定义，故\cref{definition:ch06-熵解}是适定的。

\section{持续性}\label{section:ch06-持续性}

(P) 的极弱解一个值得注意的性质是其 $2$-集中部分（见\cref{definition:ch01-集中弥散}与\cref{subsection:ch03-集中分解}）在时间上持续：

\begin{theorem}\label{theorem:ch06-6.2.1}
设 (A0)--(A1) 成立，$u$ 为问题 (P) 的极弱解。则对每个 $t\in(0,T)$，有
\begin{align}
[u(\cdot,t)]_{c,2}=[u_0]_{c,2} \quad\text{在 }\mathscr{R}_f(\Omega). \label{align:ch6-6.12}
\end{align}
\end{theorem}

\begin{remark}\label{remark:ch06-6.2.1}
设 $u$ 为 (P) 的极弱解。由\cref{theorem:ch06-6.2.1}，对所有 $t\in(0,T)$，$u(\cdot,t)$ 是 $2$-弥散的当且仅当 $u_0$ 亦然。
\end{remark}

\begin{proof}[证明\cref{theorem:ch06-6.2.1}]
设 $K\subseteq\Omega$ 为任意紧集且 $C_2(K)=0$。由\cref{proposition:ch03-零容量逼近}存在序列 $\{f_n\}\subseteq C^\infty_c(\Omega)$，满足：(i) $f_n(\Omega)\subseteq[0,1]$，$f_n|_K=1$；(ii) $\|\Delta f_n\|_{L^1(\Omega)}\to0$（$n\to\infty$）。特别，因
\begin{align*}
\int_\Omega|\nabla f_n|^2\,\mathrm{d}x=-\int_\Omega f_n\Delta f_n\,\mathrm{d}x\leqslant\int_\Omega|\Delta f_n|\,\mathrm{d}x,
\end{align*}
有 $\lim_{n\to\infty}\|f_n\|_{H^1_0(\Omega)}=0$。

设 $U\subseteq\Omega$ 为开集且 $U\supseteq K$。由\cref{lemma:appA-光滑截断}与标准正则化论证，存在 $\rho_U\in C^\infty_c(U)$，$\rho_U(U)\subseteq[0,1]$，$\rho_U|_K=1$。对所有 $t\in(0,T]$ 及 $n\in\mathbb{N}$，由等式\eqref{align:ch6-6.4}得（见\eqref{align:ch03-3.82}）
\begin{align}
\begin{aligned}
&\langle[u(\cdot,t)]_{c,2},f_n\rho_U\rangle-\langle[u_0]_{c,2},f_n\rho_U\rangle\\
&= \int_0^t\mathrm{d}s\int_\Omega\phi(u_r(x,s))\Delta(f_n\rho_U)(x)\,\mathrm{d}x\\
&\quad+\langle[u_0]_{d,2},f_n\rho_U\rangle-\langle[u(\cdot,t)]_{d,2},f_n\rho_U\rangle.
\end{aligned} \label{align:ch6-6.13}
\end{align}
因 $\|f_n\|_{H^1_0(\Omega)}\to0$，有 $f_n\rho_U\stackrel{*}{\rightharpoonup}0$ 在 $L^\infty(\Omega)$，$f_n\rho_U\to0$ 在 $H^1_0(\Omega)$。又由\cref{theorem:ch03-扩散部分H负1}，$[u_0]_{d,2}$ 与 $[u(\cdot,t)]_{d,2}$（$t\in(0,T)$）均属于 $L^1(\Omega)+H^{-1}(\Omega)$。故
\begin{align}
\lim_{n\to\infty}\langle[u(\cdot,t)]_{d,2},f_n\rho_U\rangle
=\lim_{n\to\infty}\langle[u_0]_{d,2},f_n\rho_U\rangle=0
\quad(t\in(0,T)). \label{align:ch6-6.14}
\end{align}
另一方面，因 $\phi(u_r)\in L^\infty(Q)$，$\|\Delta f_n\|_{L^1(\Omega)}\to0$，$\|f_n\|_{H^1_0(\Omega)}\to0$，得
\begin{align}
\lim_{n\to\infty}\int_0^t\int_\Omega\phi(u_r)\Delta(f_n\rho_U)\,\mathrm{d}x\,\mathrm{d}s=0. \label{align:ch6-6.15}
\end{align}
注意因对所有 $n\in\mathbb{N}$ 有 $\chi_K\leqslant f_n\rho_U\leqslant\chi_U$，故对所有 $n\in\mathbb{N}$，
\begin{align}
[u_0]^\pm_{c,2}(K)\leqslant\langle[u_0]^\pm_{c,2},f_n\rho_U\rangle\leqslant[u_0]^\pm_{c,2}(U), \label{align:ch6-6.16a}\\
[u(\cdot,t)]^\pm_{c,2}(K)\leqslant\langle[u(\cdot,t)]^\pm_{c,2},f_n\rho_U\rangle\leqslant[u(\cdot,t)]^\pm_{c,2}(U)
\quad(t\in(0,T)). \label{align:ch6-6.16b}
\end{align}
又对任意 $t\in(0,T)$，
\begin{align}
\limsup_{n\to\infty}\langle[u(\cdot,t)]_{c,2},f_n\rho_U\rangle
\leqslant\limsup_{n\to\infty}\langle[u_0]_{c,2},f_n\rho_U\rangle, \label{align:ch6-6.17}
\end{align}
因为由\eqref{align:ch6-6.14}--\eqref{align:ch6-6.15}，\eqref{align:ch6-6.13}右端当 $n\to\infty$ 时趋于零。由\eqref{align:ch6-6.16a}--\eqref{align:ch6-6.17}，对任意 $t\in(0,T)$ 得
\begin{align}
\begin{aligned}
&[u(\cdot,t)]^+_{c,2}(K)-[u(\cdot,t)]^-_{c,2}(U)
\leqslant\langle[u(\cdot,t)]_{c,2},f_n\rho_U\rangle\\
&\quad\leqslant\limsup_{n\to\infty}\langle[u(\cdot,t)]_{c,2},f_n\rho_U\rangle
\leqslant\limsup_{n\to\infty}\langle[u_0]_{c,2},f_n\rho_U\rangle\\
&\quad\leqslant\limsup_{n\to\infty}\langle[u_0]^+_{c,2},f_n\rho_U\rangle
-\liminf_{n\to\infty}\langle[u_0]^-_{c,2},f_n\rho_U\rangle
\leqslant[u_0]^+_{c,2}(U)-[u_0]^-_{c,2}(K).
\end{aligned} \label{align:ch6-6.18}
\end{align}
类似地，
\begin{align}
[u_0]^+_{c,2}(K)-[u_0]^-_{c,2}(U)
\leqslant[u(\cdot,t)]^+_{c,2}(U)-[u(\cdot,t)]^-_{c,2}(K)
\quad(t\in(0,T)). \label{align:ch6-6.19}
\end{align}
由不等式\eqref{align:ch6-6.18}--\eqref{align:ch6-6.19}及测度 $[u_0]^\pm_{c,2}$、$[u(\cdot,t)]^\pm_{c,2}$ 的正则性（见\cref{lemma:ch01-正则性}），有
\begin{align*}
[u(\cdot,t)]^+_{c,2}(K)+[u_0]^-_{c,2}(K)
\leqslant[u_0]^+_{c,2}(K)+[u(\cdot,t)]^-_{c,2}(K)
\end{align*}
及
\begin{align*}
[u_0]^+_{c,2}(K)+[u(\cdot,t)]^-_{c,2}(K)
\leqslant[u(\cdot,t)]^+_{c,2}(K)+[u_0]^-_{c,2}(K).
\end{align*}
由这些不等式，对所有 $t\in(0,T)$，有
\begin{align*}
[u(\cdot,t)]_{c,2}(K)=[u_0]_{c,2}(K)
\quad\text{对每个 }2\text{-零容紧集 }K\subseteq\Omega,
\end{align*}
从而由 $K$ 的任意性及 $[u_0]_{c,2}$、$[u(\cdot,t)]_{c,2}$ 的正则性，
\begin{align}
[u(\cdot,t)]_{c,2}(E)=[u_0]_{c,2}(E)
\quad\text{对每个 }2\text{-零容集 }E\in\mathscr{B}(\Omega). \label{align:ch6-6.20}
\end{align}
另一方面，由 $2$-集中测度的定义，对每个 $t\in(0,T)$，存在 $2$-零容集 $E_1(t)\in\mathscr{B}(\Omega)$，使 $[u(\cdot,t)]_{c,2}=[u(\cdot,t)]_{c,2}\llcorner E_1(t)$。类似地，存在 $2$-零容集 $E_2\in\mathscr{B}(\Omega)$，使 $[u_0]_{c,2}=[u_0]_{c,2}\llcorner E_2$。由等式\eqref{align:ch6-6.20}，
\begin{align*}
[u(\cdot,t)]_{c,2}(E_1(t)\setminus E_2)=[u_0]_{c,2}(E_1(t)\setminus E_2)=0,
\end{align*}
\begin{align*}
[u_0]_{c,2}(E_2\setminus E_1(t))=[u(\cdot,t)]_{c,2}(E_2\setminus E_1(t))=0.
\end{align*}
因此 $[u(\cdot,t)]_{c,2}$ 与 $[u_0]_{c,2}$ 均集中于集 $F(t):=E_1(t)\cap E_2$。显然，对任意 $E\in\mathscr{B}(\Omega)$，有 $C_2(E\cap F(t))=C_2(F(t))=0$。由\eqref{align:ch6-6.20}，
\begin{align*}
[u(\cdot,t)]_{c,2}(E)=[u(\cdot,t)]_{c,2}(E\cap F(t))
=[u_0]_{c,2}(E\cap F(t))=[u_0]_{c,2}(E)
\end{align*}
对每个 $E\in\mathscr{B}(\Omega)$ 成立。结论随之成立。
\end{proof}

\section{唯一性}\label{section:ch06-唯一性}

\cref{theorem:ch06-6.2.1}中的持续性性质在证明 (P) 弱解唯一性时起重要作用，因为它完全刻画了弱解的 $2$-集中部分。另一方面，由于\cref{definition:ch06-弱解}对 $2$-弥散部分 $[u_s(\cdot,t)]_{d,2}$ 未作任何规定（见\eqref{align:ch03-3.82}），不出所料，(P) 的弱解不构成问题的唯一类（见\cref{remark:ch06-6.3.2}）。这使如下定义变得重要。

\begin{definition}\label{definition:ch06-相容性条件}
设 (A0)--(A1) 成立。(P) 的弱解 $u$ 满足相容性条件，若对几乎所有 $t\in(0,T)$，
\begin{align}
[u^\pm_s(\cdot,t)]_{d,2}=[u^\pm_s(\cdot,t)]_{d,2}\llcorner\mathscr{S}^\pm_t, \label{align:ch6-6.21a}
\end{align}
其中
\begin{align}
\mathscr{S}^\pm_t:=\{x\in\Omega\mid\phi^\pm(u_r(x,t))=\pm\phi_{\pm\infty}\}. \label{align:ch6-6.21b}
\end{align}
\end{definition}

在本定义及此后，凡符号 $\pm$ 与 $\mp$ 出现的陈述，均应读作一对独立陈述，分别对应上、下号选择。

\begin{remark}\label{remark:ch06-6.3.1}
由\cref{definition:ch06-弱解}(i)，对几乎所有 $t\in(0,T)$，$\phi(u_r(\cdot,t))\in H^1_0(\Omega)$，故 $\phi(u_r(\cdot,t))$ 有一个 $2$-拟连续代表元 $v(\cdot,t)\in H^1_0(\Omega)$（见\cref{definition:ch02-拟连续}）。在\eqref{align:ch6-6.21b}中，我们将函数 $\phi^\pm(u_r(\cdot,t))$ 等同于其 $2$-拟连续代表元 $v^\pm(\cdot,t)\in H^1_0(\Omega)$，后者在 $\mathbb{R}$ 上 $2$-拟处处满足 $0\leqslant v^\pm\leqslant\pm\phi_{\pm\infty}$。在\eqref{align:ch6-6.21a}中，我们总将集 $\mathscr{S}^\pm_t$ 等同于 $\widetilde{\mathscr{S}}^\pm_t:=\{x\in\Omega\mid v^\pm(x,t)=\pm\phi_{\pm\infty}\}$。因 $v^\pm(\cdot,t)$ 定义到相差一个 $2$-零容集，集 $\widetilde{\mathscr{S}}^\pm_t$ 亦然。
\end{remark}

问题 (P) 的唯一性由如下定理给出，它利用\cref{definition:ch06-相容性条件}中的相容性条件。

\begin{theorem}\label{theorem:ch06-6.3.1}
设 (A0)--(A1) 成立。则对每个 $u_0\in\mathscr{R}_f(\Omega)$，问题 (P) 至多存在一个满足相容性条件\eqref{align:ch6-6.21a}的弱解。
\end{theorem}

\begin{remark}\label{remark:ch06-6.3.2}
满足相容性条件的 (P) 弱解的存在性将在下文证明（见\cref{theorem:ch06-6.4.2}）。另一方面，容易构造不满足相容性条件的 (P) 弱解，因而它们不可能是\cref{theorem:ch06-6.4.2}所给出的解。故弱解不构成问题 (P) 的唯一类。

设 (A0) 成立且 $u_{0s}\neq0$。由
\begin{align}
u(\cdot,t):=w(\cdot,t)+u_{0s} \quad\text{对几乎所有 }t\in(0,T) \label{align:ch6-6.22}
\end{align}
定义 $u\in L^\infty_{w^*}(0,T;\mathscr{R}_f(\Omega))$，其中 $w\in L^\infty(0,T;L^1(\Omega))$ 为初值 $u_{0r}$ 对应的 (P) 唯一弱解（$w$ 的存在唯一性由\cref{theorem:ch06-6.4.3}保证）。显然 $u$ 是 (P) 的弱解，$u_r=w$ 在 $Q$ 中几乎处处成立，奇异部分 $u_s(\cdot,t)=u_{0s}$ 在时间上为常数，因而 $u$ 满足相容性条件当且仅当
\begin{align*}
[u^\pm_{0s}]_{d,2}=[u^\pm_{0s}]_{d,2}\llcorner\{x\in\Omega\mid\phi^\pm(w(x,t))=\pm\phi_{\pm\infty}\}
\quad\text{对几乎所有 }t\in(0,T).
\end{align*}
这些等式不成立的例子容易给出，故断言成立。
\end{remark}

\begin{proof}[证明\cref{theorem:ch06-6.3.1}]
设 $u_1$ 与 $u_2$ 为问题 (P) 的两个弱解。为简便记 $u_{1r}\equiv(u_1)_r$，$u_{1s}\equiv(u_1)_s$（$u_2$ 类似）。设 $u_1$ 与 $u_2$ 满足相容性条件
\begin{align}
[u^\pm_{is}(\cdot,t)]_{d,2}=[u^\pm_{is}(\cdot,t)]_{d,2}\llcorner\mathscr{S}^t_{i,\pm}, \label{align:ch6-6.23}
\end{align}
其中
\begin{align*}
\mathscr{S}^t_{i,\pm}:=\{x\in\Omega\mid\phi(u_{ir})(x,t)=\phi_{\pm\infty}\}\quad(i=1,2).
\end{align*}
由\eqref{align:ch6-6.2}立即得
\begin{align}
\int_0^T\langle u_1(\cdot,t)-u_2(\cdot,t),\partial_t\zeta(\cdot,t)\rangle\,\mathrm{d}t
= \iint_Q\nabla[\phi(u_{1r})-\phi(u_{2r})]\cdot\nabla\zeta\,\mathrm{d}x\,\mathrm{d}t \label{align:ch6-6.24}
\end{align}
对每个 $\zeta\in C^1([0,T];C^1_c(\Omega))$，$\zeta(\cdot,T)=0$ 在 $\Omega$ 上成立。又由等式\eqref{align:ch6-6.12}，对所有 $t\in(0,T)$，有
\begin{align}
[u_1(\cdot,t)]_{c,2}=[u_0]_{c,2}=[u_2(\cdot,t)]_{c,2} \quad\text{在 }\mathscr{R}_f(\Omega). \label{align:ch6-6.25}
\end{align}
由\eqref{align:ch6-6.24}--\eqref{align:ch6-6.25}得
\begin{align}
\int_0^T\langle[u_1(\cdot,t)]_{d,2}-[u_2(\cdot,t)]_{d,2},\partial_t\zeta(\cdot,t)\rangle\,\mathrm{d}t
= \iint_Q\nabla[\phi(u_{1r})-\phi(u_{2r})]\cdot\nabla\zeta\,\mathrm{d}x\,\mathrm{d}t \label{align:ch6-6.26}
\end{align}
对每个 $\zeta\in C^1([0,T];C^1_c(\Omega))$，$\zeta(\cdot,T)=0$ 在 $\Omega$ 上成立。

由标准正则化论证，此等式对任意 $\zeta\in C([0,T];H^1_0(\Omega))$，$\partial_t\zeta\in L^\infty(Q)\cap L^2(0,T;H^1_0(\Omega))$，$\zeta(\cdot,T)=0$ 几乎处处在 $\Omega$ 上也成立。这使我们可在\eqref{align:ch6-6.26}中选取试验函数
\begin{align*}
\zeta(x,t):=-\int_t^T[\phi(u_{1r})-\phi(u_{2r})](x,s)\,\mathrm{d}s \quad((x,t)\in Q),
\end{align*}
从而得到
\begin{align}
\begin{aligned}
&\int_0^T\langle[u_1(\cdot,t)]_{d,2}-[u_2(\cdot,t)]_{d,2},
[\phi(u_{1r})-\phi(u_{2r})](\cdot,t)\rangle\,\mathrm{d}t\\
&\quad= -\iint_Q\nabla[\phi(u_{1r})-\phi(u_{2r})]\cdot
\left(\int_t^T\nabla[\phi(u_{1r})-\phi(u_{2r})](x,s)\,\mathrm{d}s\right)\mathrm{d}x\,\mathrm{d}t.
\end{aligned} \label{align:ch6-6.27}
\end{align}
由等式\eqref{align:ch03-3.82}有
\begin{align}
\begin{aligned}
&\int_0^T\langle[u_1(\cdot,t)]_{d,2}-[u_2(\cdot,t)]_{d,2},
[\phi(u_{1r})-\phi(u_{2r})](\cdot,t)\rangle\,\mathrm{d}t\\
&\quad= \int_0^T\langle[u_{1s}(\cdot,t)]_{d,2}-[u_2s(\cdot,t)]_{d,2},
[\phi(u_{1r})-\phi(u_{2r})](\cdot,t)\rangle\,\mathrm{d}t\\
&\qquad+\iint_Q(u_{1r}-u_{2r})[\phi(u_{1r})-\phi(u_{2r})]\,\mathrm{d}x\,\mathrm{d}t,
\end{aligned} \label{align:ch6-6.28}
\end{align}
而由相容性条件\eqref{align:ch6-6.23}，
\begin{align}
\begin{aligned}
&\int_0^T\langle[u_{1s}(\cdot,t)]_{d,2}-[u_{2s}(\cdot,t)]_{d,2},
[\phi(u_{1r})-\phi(u_{2r})](\cdot,t)\rangle\,\mathrm{d}t\\
&\quad= \int_0^T\langle[u^+_{1s}(\cdot,t)]_{d,2},\phi_{+\infty}-\phi(u_{2r})(\cdot,t)\rangle\,\mathrm{d}t\\
&\qquad-\int_0^T\langle[u^-_{1s}(\cdot,t)]_{d,2},\phi_{-\infty}-\phi(u_{2r})(\cdot,t)\rangle\,\mathrm{d}t\\
&\qquad-\int_0^T\langle[u^+_{2s}(\cdot,t)]_{d,2},\phi(u_{1r})(\cdot,t)-\phi_{+\infty}\rangle\,\mathrm{d}t\\
&\qquad+\int_0^T\langle[u^-_{2s}(\cdot,t)]_{d,2},\phi(u_{1r})(\cdot,t)-\phi_{-\infty}\rangle\,\mathrm{d}t
\geqslant0.
\end{aligned} \label{align:ch6-6.29}
\end{align}
由\eqref{align:ch6-6.28}--\eqref{align:ch6-6.29}得
\begin{align}
\begin{aligned}
&\int_0^T\langle[u_1(\cdot,t)]_{d,2}-[u_2(\cdot,t)]_{d,2},
[\phi(u_{1r})-\phi(u_{2r})](\cdot,t)\rangle\,\mathrm{d}t\\
&\quad\geqslant\iint_Q(u_{1r}-u_{2r})[\phi(u_{1r})-\phi(u_{2r})]\,\mathrm{d}x\,\mathrm{d}t.
\end{aligned} \label{align:ch6-6.30}
\end{align}
关于\eqref{align:ch6-6.27}右端，有
\begin{align}
\begin{aligned}
&-\iint_Q\nabla[\phi(u_{1r})-\phi(u_{2r})]\cdot
\left(\int_t^T\nabla[\phi(u_{1r})-\phi(u_{2r})](x,s)\,\mathrm{d}s\right)\mathrm{d}x\,\mathrm{d}t\\
&\quad= \frac{1}{2}\iint_Q\frac{\mathrm{d}}{\mathrm{d}t}
\left|\int_t^T\nabla[\phi(u_{1r})-\phi(u_{2r})](x,s)\,\mathrm{d}s\right|^2\mathrm{d}x\,\mathrm{d}t\\
&\quad= -\frac{1}{2}\int_\Omega\left|\int_0^T\nabla[\phi(u_{1r})-\phi(u_{2r})](x,s)\,\mathrm{d}s\right|^2\mathrm{d}x.
\end{aligned} \label{align:ch6-6.31}
\end{align}
由\eqref{align:ch6-6.27}与\eqref{align:ch6-6.30}--\eqref{align:ch6-6.31}得
\begin{align*}
\iint_Q(u_{1r}-u_{2r})[\phi(u_{1r})-\phi(u_{2r})]\,\mathrm{d}x\,\mathrm{d}t
\leqslant-\frac{1}{2}\int_\Omega\left|\int_0^T\nabla[\phi(u_{1r})-\phi(u_{2r})](x,t)\,\mathrm{d}t\right|^2\mathrm{d}x\leqslant0,
\end{align*}
从而
\begin{align}
\phi(u_{1r}(x,t))=\phi(u_{2r}(x,t)) \quad\text{对几乎所有 }(x,t)\in Q. \label{align:ch6-6.32}
\end{align}
现在在等式\eqref{align:ch6-6.24}中取 $\zeta(x,t)=\rho(x)k_q(t)$，其中 $\rho\in C^2_c(\Omega)$，$k_q$ 如\eqref{align:ch6-6.7}。如\eqref{align:ch6-6.4}证明中那样论证，并利用等式\eqref{align:ch6-6.32}，对几乎所有 $t_1\in(0,T)$，得
\begin{align*}
\langle u_1(\cdot,t_1)-u_2(\cdot,t_1),\rho\rangle=0.
\end{align*}
由 $\rho$ 的任意性，这蕴含对几乎所有 $t_1\in(0,T)$，$u_1(\cdot,t_1)=u_2(\cdot,t_1)$ 在 $\mathscr{R}_f(\Omega)$。结论随之成立。
\end{proof}

\section{存在性与正则性结果}\label{section:ch06-存在正则}

\subsection{存在性}\label{subsection:ch06-存在性}

第一个存在性结果的证明是构造性的。事实上，我们通过研究逼近问题序列的解 $\{u_n\}$ 来展示 (P) 的弱解：
\begin{align*}
\begin{cases}
\partial_t u_n=\Delta\phi_n(u_n) & \text{在 }Q,\\
u_n=0 & \text{在 }\Gamma,\\
u_n(\cdot,0)=u_{0n} & \text{在 }\Omega,
\end{cases}
\qquad(Pn)
\end{align*}
（见\cref{section:ch06-逼近问题}）。此处 $\phi_n(z):=\phi(z)+\frac{z}{n}$（$z\in\mathbb{R}$），$u_{0n}$ 为初值测度 $u_0$ 的合适正则化（见\eqref{align:ch6-6.48}--\eqref{align:ch6-6.49a}）。

\begin{theorem}\label{theorem:ch06-6.4.1}
设 (A0)--(A2) 成立。则存在问题 (P) 的弱解，它作为逼近问题 (Pn) 解序列在 $L^\infty_{w^*}(0,T;\mathscr{R}_f(\Omega))$ 的弱$^*$拓扑下的极限点而得到。进而：
\begin{enumerate}[(i)]
\item 对几乎所有 $0<t_1<t_2<T$，
\begin{align}
u^\pm_s(\cdot,t_2)\leqslant u^\pm_s(\cdot,t_1)\leqslant u^\pm_{0s} \quad\text{在 }\mathscr{R}_f(\Omega). \label{align:ch6-6.33}
\end{align}
\item 若 $u_0\in L^1(\Omega)$，则 $u\in L^\infty(0,T;L^1(\Omega))$。
\end{enumerate}
\end{theorem}

由\cref{theorem:ch06-6.3.1}，为得到 (P) 的适定性，证明满足相容性条件的弱解存在性是重要的。这由如下结果保证。

\begin{theorem}\label{theorem:ch06-6.4.2}\label{proposition:ch06-6.4.2}
设 (A0)--(A2) 成立。则\cref{theorem:ch06-6.4.1}所给出的问题 (P) 的弱解满足相容性条件。
\end{theorem}

由\cref{theorem:ch06-6.3.1}、\cref{theorem:ch06-6.4.1}与\cref{theorem:ch06-6.4.2}，有如下结果。

\begin{theorem}\label{theorem:ch06-6.4.3}
设 (A0)--(A2) 成立。则对每个 $u_0\in\mathscr{R}_f(\Omega)$，存在问题 (P) 满足相容性条件的唯一弱解 $u$。进而\cref{theorem:ch06-6.4.1}的 (i)--(ii) 款成立。
\end{theorem}

\cref{theorem:ch06-6.4.1}证明中构造的 (P) 弱解有若干重要性质（特别地，由\cref{theorem:ch06-6.4.2}它们满足相容性条件）。这引出如下定义。

\begin{definition}\label{definition:ch06-构造解}
\cref{theorem:ch06-6.4.1}所给出的问题 (P) 弱解称为构造解。
\end{definition}

构造解的另一个有趣性质如下。

\begin{proposition}\label{proposition:ch06-6.4.4}
设 (A0)--(A2) 成立。则问题 (P) 的每个构造解都是熵解。
\end{proposition}

\subsection{正则性}\label{subsection:ch06-正则性}

熵解一个有趣的特征是其奇异部分既不能自发出现，也不能随时间增长：

\begin{theorem}\label{theorem:ch06-6.4.5}
设 (A0)--(A1) 成立。假设
\begin{align}
\text{存在 }L_0\geqslant0\text{ 使当 }|s|>L_0\text{ 时 }\phi'(s)>0. \label{align:ch6-6.34}
\end{align}
设 $u$ 为问题 (P) 的熵解。则对几乎所有 $t_1,t_2\in(0,T)$，$t_1<t_2$，不等式\eqref{align:ch6-6.33}成立。
\end{theorem}

\begin{proof}
设 $u$ 为问题 (P) 的熵解。先证对所有 $\rho\in C^1_c(\Omega)$，$\rho\geqslant0$，及 $h\in C^1([0,T])$，$h(T)=0$，有
\begin{align}
\int_0^T\langle u^\pm_s(\cdot,t),\rho\rangle h'(t)\,\mathrm{d}t\geqslant-h(0)\langle u^\pm_{0s},\rho\rangle. \label{align:ch6-6.35}
\end{align}
我们仅证 $+$ 号情形。因 $\phi$ 在 $[L_0,+\infty)$ 上严格增（见\eqref{align:ch6-6.34}），由标准逼近结果可在\eqref{align:ch6-6.11}中取 $F_\phi(y)=[y-m]_+$，$m>L_0$（这对应于在\eqref{align:ch6-6.9}中取 $f=\chi_{[\phi(m),+\infty)}$）。此时 $M^+_{F_\phi}=1$，$M^-_{F_\phi}=0$，从而得
\begin{align}
\begin{aligned}
&\iint_Q\{[u_r-m]_+\zeta_t-\chi_{[m,+\infty)}(u_r)\nabla\phi(u_r)\cdot\nabla\zeta\}\,\mathrm{d}x\,\mathrm{d}t\\
&\quad+\int_0^T\langle u^+_s(\cdot,t),\zeta_t(\cdot,t)\rangle\,\mathrm{d}t
\geqslant-\int_\Omega[u_{0r}-m]_+\zeta(x,0)\,\mathrm{d}x-\langle u^+_{0s},\zeta(\cdot,0)\rangle
\end{aligned} \label{align:ch6-6.36}
\end{align}
对每个 $\zeta\in C^1([0,T];C^1_c(\Omega))$，$\zeta\geqslant0$，$\zeta(\cdot,T)=0$ 在 $\Omega$ 上成立。在此不等式中取 $\zeta(x,t)=\rho(x)h(t)$，$\rho,h$ 如上，并令 $m\to\infty$，即得\eqref{align:ch6-6.35}。

现在固定 $t_1\in(0,T)$。对充分大的 $p\in\mathbb{N}$，置
\begin{align*}
h_p(t):=\chi_{[0,t_1]}(t)+p\left(t_1+\frac{1}{p}-t\right)\chi_{(t_1,t_1+\frac{1}{p}]}(t) \quad(t\in[0,T]).
\end{align*}
在\eqref{align:ch6-6.35}中取 $h=h_p$，对所有非负 $\rho\in C^1_c(\Omega)$，得
\begin{align}
p\int_{t_1}^{t_1+\frac{1}{p}}\langle u^\pm_s(\cdot,t),\rho\rangle\,\mathrm{d}t\leqslant\langle u^\pm_{0s},\rho\rangle. \label{align:ch6-6.37}
\end{align}
因由\cref{definition:ch06-熵解}，映射 $t\mapsto\langle u^\pm_s(\cdot,t),\rho\rangle\in L^\infty(0,T)$，由\cref{corollary:ch02-平均收敛}(i) 有
\begin{align*}
\lim_{p\to\infty}p\int_{t_1}^{t_1+\frac{1}{p}}\langle u^\pm_s(\cdot,t),\rho\rangle\,\mathrm{d}t
=\langle u^\pm_s(\cdot,t_1),\rho\rangle.
\end{align*}
在\eqref{align:ch6-6.37}中令 $p\to\infty$，对几乎所有 $t_1\in(0,T)$，得 $\langle u^\pm_s(\cdot,t_1),\rho\rangle\leqslant\langle u^\pm_{0s},\rho\rangle$，对所有 $\rho\in C^1_c(\Omega)$，$\rho\geqslant0$ 成立，故\eqref{align:ch6-6.33}中第二个不等式直接成立。第一个不等式类似地在\eqref{align:ch6-6.35}中取如下 $h=h_p$ 得证：
\begin{align*}
h_p(t):=p(t-t_1)\chi_{[t_1,t_1+\frac{1}{p}]}(t)
+\chi_{(t_1+\frac{1}{p},t_2)}(t)
+p\left(t_2+\frac{1}{p}-t\right)\chi_{[t_2,t_2+\frac{1}{p}]}(t)
\quad(t\in[0,T]).
\end{align*}
结论随之成立。
\end{proof}

\cref{theorem:ch06-6.4.5}可视为熵解的一个正则性结果。特别地，由\cref{proposition:ch06-6.4.4}，它适用于构造解。值得注意的是，(P) 的构造解还表现出进一步的正则化效应。为讨论这一点，陈述如下定义。

\begin{definition}\label{definition:ch06-RfL1正则}
问题 (P) 的弱解 $u$ 经历 $\mathscr{R}_f$--$L^1$ 正则化效应，若存在等待时间 $\tau\in[0,T)$，使对所有 $t\in(\tau,T)$ 有 $u(\cdot,t)\in L^1(\Omega)$。若 $\tau=0$，则发生瞬时 $\mathscr{R}_f$--$L^1$ 正则化。
\end{definition}

$\mathscr{R}_f$--$L^1$ 正则化效应发生的充分条件由如下命题给出，它依赖于 $\phi$ 趋于饱和值 $\phi_{\pm\infty}$ 的合适收敛速率（证明见文献~\cite{ref87} 命题 3.8）。

\begin{proposition}\label{proposition:ch6-瞬时正则}
设 (A0)--(A2) 成立。设存在 $s_\pm\in(0,1)$ 与 $M_\pm\in(0,\phi_{\pm\infty}]$，使
\begin{align}
\pm\phi_{\pm\infty}\mp\phi(y)\geqslant\frac{M_\pm}{(1+|y|)^{s_\pm}}
\quad\text{当 }\pm y\geqslant0. \label{align:ch6-6.38}
\end{align}
进而设 $u^\pm_0\in\mathscr{R}^+_{d,p}(\Omega)$，$p\in(1,\frac{2}{1+s_\pm})$，$u$ 为 (P) 满足相容性条件的唯一弱解。则 $u^\pm\in L^\infty(0,T;L^1(\Omega))$。
\end{proposition}

相反地，如下两个命题指出瞬时 $\mathscr{R}_f$--$L^1$ 正则化不发生的情形：或对合适的 $u_0\in\mathscr{R}_{d,2}(\Omega)$ 且 $\phi$ 严格增，或在 $\phi$ 在 $\pm\infty$ 附近为常数的最奇异情形下对任意构造解。证明见文献~\cite{ref87} 命题 3.9、3.10。

\begin{proposition}\label{proposition:ch06-6.4.7}
设 (A0)--(A2) 成立且 $\phi$ 严格增。则存在 $u_0\in\mathscr{R}^+_{d,2}(\Omega)$，$u_{0s}\neq0$，使对 (P) 满足相容性条件的唯一弱解，有
\begin{align}
\operatorname*{ess\,lim}_{t\to0^+}\|u_s(\cdot,t)-u_{0s}\|_{\mathscr{R}_f(\Omega)}=0. \label{align:ch6-6.39}
\end{align}
\end{proposition}

\begin{proposition}\label{proposition:ch06-6.4.8}
设 (A0)--(A2) 成立，$u$ 为 (P) 满足相容性条件的唯一弱解。设存在 $\pm c_\pm\geqslant0$ 与 $\pm d_\pm\geqslant0$，使 $\phi(y)=c_\pm$ 当 $\pm y\geqslant\pm d_\pm$。则对每个开子集 $\widetilde\Omega\Subset\Omega$，有
\begin{align}
\operatorname*{ess\,lim}_{t\to0^+}\|u^\pm_s(\cdot,t)-u^\pm_{0s}\|_{\mathscr{R}_f(\widetilde\Omega)}=0. \label{align:ch6-6.40}
\end{align}
\end{proposition}

\section{存在性结果的证明：逼近问题 (Pn)}\label{section:ch06-逼近问题}

\subsection{初值的逼近}\label{subsection:ch06-初值逼近}

设 (A0) 成立。为处理问题 (P)，需要初值 $u_0$ 的合适逼近。为此，使用如下技术引理。

\begin{lemma}\label{lemma:ch06-6.5.1}
对任意 $\mu\in\mathscr{R}^+_f(\Omega)$，存在序列 $\{(\alpha_n,\beta_n,\gamma_n)\}\subseteq[C^\infty_c(\Omega)]^3$，$\alpha_n\geqslant0$，$\beta_n\geqslant0$，$\gamma_n\geqslant0$ 在 $\Omega$ 上，具有如下性质：
\begin{align}
\alpha_n\to\mu_r \text{ 在 }L^1(\Omega), \label{align:ch6-6.41a}
\end{align}
\begin{align}
\alpha_n\stackrel{*}{\rightharpoonup}\mu_{\mathrm{ac}},\quad
\beta_n\stackrel{*}{\rightharpoonup}[\mu_s]_{d,2},\quad
\gamma_n\stackrel{*}{\rightharpoonup}\mu_{c,2}
\quad\text{在 }\mathscr{R}_f(\Omega), \label{align:ch6-6.41b}
\end{align}
\begin{align}
\int_\Omega\beta_n f_n\rho\,\mathrm{d}x\to\langle[\mu_s]_{d,2},f\rho\rangle \label{align:ch6-6.41c}
\end{align}
对所有 $\rho\in C_c(\Omega)$ 及所有 $\{f_n\}\subseteq H^1(\Omega)\cap L^\infty(\Omega)$，$\sup_{n\in\mathbb{N}}\|f_n\|_\infty<\infty$，$f_n\rightharpoonup f$ 在 $H^1(\Omega)$ 成立。

置 $\mu_n:=\alpha_n+\beta_n+\gamma_n\in C^\infty_c(\Omega)$。可取 $\alpha_n,\beta_n,\gamma_n$ 使
\begin{align}
\|\mu_n\|_{L^1(\Omega)}\leqslant\|\mu\|_{\mathscr{R}_f(\Omega)},\quad
\|\mu_n\|_{L^\infty(\Omega)}\leqslant\frac{M\|\mu\|_{\mathscr{R}_f(\Omega)}}{\sqrt[4]{n}}
\quad\text{对某个 }M>0, \label{align:ch6-6.41d}
\end{align}
\begin{align}
\mu_n\to\mu_r \text{ 在 }\Omega\text{ 中几乎处处},\quad
\mu_n\stackrel{*}{\rightharpoonup}\mu \text{ 在 }\mathscr{R}_f(\Omega). \label{align:ch6-6.41e}
\end{align}
\end{lemma}

\begin{proof}
置 $\widetilde\mu:=\widetilde\mu_r+\widetilde\mu_s\in\mathscr{R}^+_f(\mathbb{R}^N)$，
\begin{align*}
\widetilde\mu_r(x):=\mu_r(x)\chi_\Omega(x)\quad(x\in\mathbb{R}^N),\qquad
\widetilde\mu_s(E):=\mu_s(\Omega\cap E)\quad(E\in\mathscr{B}^d).
\end{align*}
由定义 $\widetilde\mu=\widetilde\mu\llcorner\Omega$，故对每个 $E\in\mathscr{B}^d\cap\Omega$，$\widetilde\mu(E)=\mu(E)$。考虑序列 $\{\zeta_n\}\subseteq C^\infty_c(\mathbb{R}^N)$，定义为
\begin{align}
\zeta_n(x):=\frac{n^{d\theta}}{\int_{\mathbb{R}^N}\zeta(x)\,\mathrm{d}x}\zeta(n^\theta x) \quad(x\in\mathbb{R}^N), \label{align:ch6-6.42}
\end{align}
其中 $\zeta\in C^\infty_c(\mathbb{R}^N)$ 为标准磨光子，$\theta>0$ 待选。对任意 $n\in\mathbb{N}$，置
\begin{align*}
\widetilde\alpha_n:=\widetilde\mu_{\mathrm{ac}}*\zeta_n,\quad
\widetilde\beta_n:=[\widetilde\mu_s]_{d,2}*\zeta_n,\quad
\widetilde\gamma_n:=\widetilde\mu_{c,2}*\zeta_n,\quad
\widetilde\mu_n:=\widetilde\alpha_n+\widetilde\beta_n+\widetilde\gamma_n=\widetilde\mu*\zeta_n
\end{align*}
（见\eqref{align:ch03-3.82}，$p=2$，$m=1$）。设 $\Omega_n$ 为开集，$\Omega_n\Subset\Omega_{n+1}\Subset\Omega$ 对每个 $n\in\mathbb{N}$，$\bigcup_{n=1}^\infty\Omega_n=\Omega$，设 $\eta_n\in C^\infty_c(\Omega_{n+1})$，$0\leqslant\eta_n\leqslant1$，$\eta_n=1$ 在 $\Omega_n$ 上（$n\in\mathbb{N}$）。

我们证明函数
\begin{align*}
\alpha_n:=\widetilde\alpha_n\eta_n,\quad
\beta_n:=\widetilde\beta_n\eta_n,\quad
\gamma_n:=\widetilde\gamma_n\eta_n,\quad
\mu_n:=\widetilde\mu_n\eta_n=\alpha_n+\beta_n+\gamma_n\quad(n\in\mathbb{N})
\end{align*}
（它们显然在 $\Omega$ 上非负）具有\eqref{align:ch6-6.41a}中所述性质。

\eqref{align:ch6-6.41a}--\eqref{align:ch6-6.41b}中的收敛断言由标准卷积结果得到。为证\eqref{align:ch6-6.41c}，固定 $\rho\in C^1_c(\Omega)$。置 $F_n:=(f_n\rho)*\zeta_n$，取 $n$ 充分大使 $\operatorname{supp}\rho\subseteq\Omega_n$ 且 $\operatorname{supp}F_n\subseteq\Omega$。注意 $\eta_n=1$ 在 $\Omega_n$ 上，得
\begin{align}
\langle[\mu_s]_{d,2},F_n\rangle
=\langle[\widetilde\mu_s]_{d,2}*\zeta_n,\eta_n f_n\rho\rangle
=\int_\Omega\beta_n f_n\rho\,\mathrm{d}x. \label{align:ch6-6.43}
\end{align}
注意由 Banach--Alaoglu 定理 $f_n\stackrel{*}{\rightharpoonup}f$ 在 $L^\infty(\Omega)$，因为由假设 $\{f_n\}$ 在 $L^\infty(\Omega)$ 中有界。显然 $F_n\stackrel{*}{\rightharpoonup}f\rho$ 在 $L^\infty(\Omega)$。事实上，对每个 $\zeta\in L^1(\Omega)$，有
\begin{align}
\begin{aligned}
\lim_{n\to\infty}\int_\Omega F_n\zeta\,\mathrm{d}x
&=\lim_{n\to\infty}\int_\Omega(f_n\rho)*\zeta_n\,\zeta\,\mathrm{d}x\\
&=\lim_{n\to\infty}\int_\Omega\zeta(x)\left(\int_\Omega f_n(y)\rho(y)\zeta_n(x-y)\,\mathrm{d}y\right)\mathrm{d}x\\
&=\lim_{n\to\infty}\int_\Omega f_n(y)\rho(y)\,\mathrm{d}y\int_\Omega\zeta(x)\zeta_n(x-y)\,\mathrm{d}x
=\int_\Omega f\rho\zeta\,\mathrm{d}x.
\end{aligned} \label{align:ch6-6.44}
\end{align}
由\eqref{align:ch6-6.44}也有 $F_n\rightharpoonup f\rho$ 在 $L^2(\Omega)$。因 $f_n\rho\in H^1_0(\Omega)$，对每个充分大的 $n$，$F_n\in H^1_0(\Omega)$，弱梯度 $\nabla F_n=[\nabla(f_n\rho)]*\zeta_n$。于是显然 $F_n\rightharpoonup f\rho$ 在 $H^1_0(\Omega)$，因为由假设 $f_n\rightharpoonup f$ 在 $H^1(\Omega)$。

因 $[\mu_s]_{d,2}\in L^1(\Omega)+H^{-1}(\Omega)$（见\cref{theorem:ch03-扩散部分H负1}），由上述收敛结果得
\begin{align*}
\lim_{n\to\infty}\langle[\mu_s]_{d,2},F_n\rangle=\langle[\mu_s]_{d,2},f\rho\rangle.
\end{align*}
在\eqref{align:ch6-6.43}中令 $n\to\infty$，即得\eqref{align:ch6-6.41c}对所有 $\rho\in C^1_c(\Omega)$，由标准逼近结果断言随之成立。

\eqref{align:ch6-6.41d}中第一个不等式由 $\mu_n$ 的定义直接得到，第二个通过在\eqref{align:ch6-6.42}中取 $\theta\in(0,1/(4N)]$ 得到：
\begin{align*}
\|\mu_n\|_{L^\infty(\Omega)}\leqslant Cn^{N\theta}\|\mu\|_{\mathscr{R}_f(\Omega)}
\leqslant\frac{C}{\sqrt[4]{n}}\|\mu\|_{\mathscr{R}_f(\Omega)}.
\end{align*}

现在处理\eqref{align:ch6-6.41e}。置
\begin{align}
[\mu_s]_n:=\beta_n+\gamma_n\equiv(\widetilde\mu_s*\zeta_n)\eta_n, \label{align:ch6-6.45}
\end{align}
注意由 Radon--Nikodým 定理，对几乎所有 $x\in\Omega$，
\begin{align}
\lim_{n\to\infty}\alpha_n(x)=\mu_r(x), \label{align:ch6-6.46}
\end{align}
\begin{align}
\limsup_{n\to\infty}[\mu_s]_n(x)
\leqslant\frac{1}{\int_{\mathbb{R}^N}\zeta(x)\,\mathrm{d}x}
\lim_{n\to\infty}n^{N\theta}\widetilde\mu_s\left(B\left(x,\frac{1}{n^\theta}\right)\right)=0 \label{align:ch6-6.47}
\end{align}
（见\eqref{align:ch02-2.77}中第二个不等式，\cref{corollary:ch02-Lebesgue微分}）。因
\begin{align*}
\mu_n=\alpha_n+[\mu_s]_n,
\end{align*}
结合\eqref{align:ch6-6.46}与\eqref{align:ch6-6.47}得
\begin{align*}
\mu_n(x)\to\mu_r(x) \quad\text{对几乎所有 }x\in\Omega,
\end{align*}
\eqref{align:ch6-6.41e}中第一个收敛得证（\eqref{align:ch6-6.41e}中第二个收敛另见\eqref{align:ch6-6.41b}）。

对任意 $u_0\in\mathscr{R}_f(\Omega)$，考虑 Jordan 分解 $u_0=u^+_0-u^-_0$，其中 $u^\pm_0\in\mathscr{R}^+_f(\Omega)$，$u^+_0\perp u^-_0$。设 $u^\pm_{0r}\in L^1(\Omega)$ 为 $[u^\pm_0]_{\mathrm{ac}}$ 的密度，$\delta^\pm:=[u^\pm_0]_{d,2}$，$\gamma^\pm:=[u^\pm_0]_{c,2}$ 分别为 $u^\pm_0$ 关于 $2$-容量的弥散部分与集中部分。由\eqref{align:ch03-3.82}有 $u^\pm_0=\gamma^\pm+\delta^\pm$。置
\begin{align}
[u^\pm_0]_n:=\gamma^\pm_n+\delta^\pm_n\geqslant0,\quad
u_{0n}:=[u^+_0]_n-[u^-_0]_n, \label{align:ch6-6.48}
\end{align}
其中 $\{\gamma^\pm_n\}\subseteq C^\infty_c(\Omega)$，$\{\delta^\pm_n\}\subseteq C^\infty_c(\Omega)$，$\gamma^\pm_n\geqslant0$，$\delta^\pm_n\geqslant0$，满足
\begin{align}
\|[u^\pm_0]_n\|_{L^1(\Omega)}\leqslant\|u^\pm_0\|_{\mathscr{M}(\Omega)},\quad
\|[u^\pm_0]_n\|_{L^\infty(\Omega)}\leqslant\frac{M\|u^\pm_0\|_{\mathscr{M}(\Omega)}}{\sqrt[4]{n}}, \label{align:ch6-6.49a}
\end{align}
\begin{align}
[u^\pm_0]_n\to u^\pm_{0r} \text{ 在 }\Omega\text{ 中几乎处处}, \label{align:ch6-6.49b}
\end{align}
\begin{align}
[u^\pm_0]_n\stackrel{*}{\rightharpoonup}u^\pm_0 \text{ 在 }\mathscr{M}(\Omega), \label{align:ch6-6.49c}
\end{align}
\begin{align}
\delta^\pm_n\stackrel{*}{\rightharpoonup}[u^\pm_0]_{d,2},\quad
\gamma^\pm_n\stackrel{*}{\rightharpoonup}[u^\pm_0]_{c,2}
\quad\text{在 }\mathscr{M}(\Omega), \label{align:ch6-6.49d}
\end{align}
\begin{align}
\int_\Omega\delta^\pm_n f_n\rho\,\mathrm{d}x\to\langle[u^\pm_0]_{d,2},f\rho\rangle \label{align:ch6-6.49e}
\end{align}
对所有 $\rho\in C_c(\Omega)$ 及 $\{f_n\}\subseteq H^1(\Omega)\cap L^\infty(\Omega)$，$\sup_{n\in\mathbb{N}}\|f_n\|_{L^\infty(\Omega)}<\infty$，$f_n\rightharpoonup f$ 在 $H^1(\Omega)$ 成立。此后总假设 (Pn) 的初值 $u_{0n}$ 具有性质\eqref{align:ch6-6.49a}。
\end{proof}

\begin{remark}\label{remark:ch06-6.5.1}
对任意 $u_0\in\mathscr{R}_f(\Omega)$，考虑分解（见\eqref{align:ch03-3.82}，$p=2$，$m=1$）
\begin{align}
u^\pm_0=[u^\pm_0]_{\mathrm{ac}}+[u^\pm_{0s}]_{d,2}+[u^\pm_0]_{c,2}. \label{align:ch6-6.50}
\end{align}
对 $\mu=u^\pm_0$ 应用\cref{lemma:ch06-6.5.1}，得序列 $\{(\alpha^\pm_n,\beta^\pm_n,\gamma^\pm_n)\}\subseteq[C^\infty_c(\Omega)]^3$ 满足\eqref{align:ch6-6.41a}。置 $\delta^\pm_n:=\alpha^\pm_n+\beta^\pm_n$，性质\eqref{align:ch6-6.49a}立得。
\end{remark}

\begin{lemma}\label{lemma:ch06-6.5.2}
设 $\{u_{0n}\}$ 满足\eqref{align:ch6-6.48}与\eqref{align:ch6-6.49a}--\eqref{align:ch6-6.49c}。设 $f\in C(\mathbb{R})$ 满足\cref{proposition:ch05-函数弱星}。则存在 $\{u_{0n}\}$ 的子列（不重新标号），使
\begin{align}
f\circ u_{0n}\stackrel{*}{\rightharpoonup}f\circ u_{0r}+M^+_f u^+_{0s}-M^-_f u^-_{0s}
\quad\text{在 }\mathscr{R}_f(\Omega). \label{align:ch6-6.51}
\end{align}
特别地，$u_{0n}\stackrel{*}{\rightharpoonup}u_0$ 在 $\mathscr{R}_f(\Omega)$。
\end{lemma}

\begin{proof}
设 $\{\nu_n\}\subseteq\mathcal{Y}^+(\Omega;\mathbb{R})$ 为与 $\{u_{0n}\}$ 相伴的 Young 测度序列。由\eqref{align:ch6-6.49a}中第一个不等式，序列 $\{u_{0n}\}$ 在 $L^1(\Omega)$ 中有界，故由\cref{theorem:ch05-弱紧Young}存在 Young 测度 $\nu\in\mathcal{Y}^+(\Omega;\mathbb{R})$，使（可能抽取子列，不重新标号）$\nu_n\to\nu$ 窄收敛。由\eqref{align:ch6-6.49b}中的收敛、\cref{proposition:ch05-几乎处处收敛}及等式\eqref{align:ch05-5.53}，$\nu$ 的离散分解 $\{\nu_x\}_{x\in\Omega}$ 满足对几乎所有 $x\in\Omega$，$\nu_x=\delta_{u_{0r}(x)}$。于是
\begin{align*}
\int_\mathbb{R}y^\pm\,\mathrm{d}\nu_x(y)=u^\pm_{0r}(x)
\quad\text{对几乎所有 }x\in\Omega.
\end{align*}
因此，如\cref{remark:ch05-biting注}中那样论证，存在 $\{u_{0n}\}$ 的子列（不重新标号）及 $\sigma^\pm\in\mathscr{R}^+_f(\Omega)$，使
\begin{align}
[u_{0n}]^\pm\stackrel{*}{\rightharpoonup}u^\pm_{0r}+\sigma^\pm \quad\text{在 }\mathscr{R}_f(\Omega). \label{align:ch6-6.52}
\end{align}
由 $u_{0n}$ 的定义（见\eqref{align:ch6-6.48}）及\eqref{align:ch6-6.49c}有 $u_{0n}\stackrel{*}{\rightharpoonup}u^+_0-u^-_0=u_0$，而由\eqref{align:ch6-6.52}，$u_{0n}\stackrel{*}{\rightharpoonup}u_{0r}+\sigma$，其中 $\sigma:=\sigma^+-\sigma^-$。

注意由\eqref{align:ch6-6.48}中第二个等式，对所有 $n\in\mathbb{N}$ 有 $[u_{0n}]^\pm\leqslant[u^\pm_0]_n$，故由\eqref{align:ch6-6.49b}--\eqref{align:ch6-6.49c}与\eqref{align:ch6-6.52}得 $\sigma^\pm\leqslant u^\pm_{0s}$ 在 $\mathscr{R}_f(\Omega)$。这蕴含 $\sigma^\pm$ 关于 Lebesgue 测度奇异且相互奇异。于是等式 $u_0=u_{0r}+\sigma$ 及 Lebesgue 分解的唯一性给出 $\sigma=u_{0s}$，$\sigma^\pm=u^\pm_{0s}$。然后如\cref{proposition:ch05-集中测度极限}证明中那样，结论随之成立。
\end{proof}

\subsection{先验估计}\label{subsection:ch06-先验估计}

由经典结果，对任意 $n\in\mathbb{N}$，问题 (Pn) 存在唯一经典解 $u_n\in C^\infty(Q)$（例如见文献~\cite{ref63}）；且对任意 $n\in\mathbb{N}$，有
\begin{align}
\|u_n\|_{L^\infty(0,T;L^1(\Omega))}\leqslant\|u_{0n}\|_{L^1(\Omega)},\quad
\|u_n\|_{L^\infty(Q)}\leqslant\|u_{0n}\|_{L^\infty(\Omega)}. \label{align:ch6-6.53}
\end{align}
我们需要序列 $\{u_n\}$ 与 $\{\phi(u_n)\}$ 的先验估计。这由如下命题给出。

\begin{lemma}\label{lemma:ch06-6.5.3}
设 (A0)--(A1) 成立，$\{u_n\}$ 为问题 (Pn) 的解序列。则对所有 $n\in\mathbb{N}$，有
\begin{align}
\|u_n\|_{L^\infty(0,T;L^1(\Omega))}\leqslant\|u_0\|_{\mathscr{R}_f(\Omega)}, \label{align:ch6-6.54}
\end{align}
\begin{align}
\|u_n\|_{L^\infty(Q)}\leqslant\frac{M\|u_0\|_{\mathscr{R}_f(\Omega)}}{\sqrt[4]{n}}, \label{align:ch6-6.55}
\end{align}
从而
\begin{align}
\|\phi_n(u_n)-\phi(u_n)\|_{L^\infty(Q)}
\leqslant\frac{\|u_n\|_{L^\infty(Q)}}{n}
\leqslant M\|u_0\|_{\mathscr{R}_f(\Omega)}n^{-3/4}, \label{align:ch6-6.56}
\end{align}
\begin{align}
\iint_Q|\nabla[\phi(u_n)]|^2\,\mathrm{d}x\,\mathrm{d}t
+\frac{1}{n}\iint_Q\phi'(u_n)|\nabla u_n|^2\,\mathrm{d}x\,\mathrm{d}t
\leqslant2\|\phi\|_\infty\|u_0\|_{\mathscr{R}_f(\Omega)}. \label{align:ch6-6.57}
\end{align}
\end{lemma}

\begin{proof}
不等式\eqref{align:ch6-6.49a}与\eqref{align:ch6-6.53}立即给出\eqref{align:ch6-6.54}与\eqref{align:ch6-6.55}，\eqref{align:ch6-6.56}由 $\phi_n$ 的定义直接得到（回忆 $\phi_n(z):=\phi(z)+z/n$，$z\in\mathbb{R}$）。为证\eqref{align:ch6-6.57}，将 (Pn) 第一式乘以 $\phi(u_n)$ 并在 $Q$ 上积分。利用\eqref{align:ch6-6.54}，得
\begin{align*}
\begin{aligned}
&\iint_Q|\nabla[\phi(u_n)]|^2\,\mathrm{d}x\,\mathrm{d}t
+\frac{1}{n}\iint_Q\phi'(u_n)|\nabla u_n|^2\,\mathrm{d}x\,\mathrm{d}t\\
&\quad=-\iint_Q\phi(u_n)\partial_t u_n\,\mathrm{d}x\,\mathrm{d}t\\
&\quad=\int_\Omega\mathrm{d}x\int_0^{u_{0n}(x)}\phi(y)\,\mathrm{d}y
-\int_\Omega\mathrm{d}x\int_0^{u_n(x,T)}\phi(y)\,\mathrm{d}y\\
&\quad\leqslant\|\phi\|_\infty[\|u_{0n}\|_{L^1(\Omega)}+\|u_n(\cdot,T)\|_{L^1(\Omega)}]
\leqslant2\|\phi\|_\infty\|u_0\|_{\mathscr{R}_f(\Omega)}.
\end{aligned}
\end{align*}
结论随之成立。
\end{proof}

对每个 $h\in C(\mathbb{R})$，置
\begin{align}
H(y):=\int_0^y h(z)\,\mathrm{d}z \quad(y\in\mathbb{R}). \label{align:ch6-6.58}
\end{align}

\begin{lemma}\label{lemma:ch06-6.5.4}
设 (A1) 成立，$\{u_n\}$ 为问题 (Pn) 的解序列。
\begin{enumerate}[(i)]
\item 对所有 $h\in\operatorname{Lip}(\mathbb{R})$，$\zeta\in C^1([0,T];C^1_c(\Omega))$，及 $t\in(0,T]$，有
\begin{align}
\begin{aligned}
&\int_\Omega H(u_n(x,t))\zeta(x,t)\,\mathrm{d}x
+\iint_{Q_t}h'(u_n)\nabla u_n\cdot\nabla\phi_n(u_n)\zeta\,\mathrm{d}x\,\mathrm{d}s\\
&\qquad-\iint_{Q_t}H(u_n)\partial_s\zeta\,\mathrm{d}x\,\mathrm{d}s\\
&\quad=\int_\Omega H(u_{0n}(x))\zeta(x,0)\,\mathrm{d}x
-\iint_{Q_t}h(u_n)\nabla\phi_n(u_n)\cdot\nabla\zeta\,\mathrm{d}x\,\mathrm{d}s,
\end{aligned} \label{align:ch6-6.59}
\end{align}
其中 $Q_t:=\Omega\times(0,t)$。若 $h(0)=0$，则\eqref{align:ch6-6.59}对所有 $\zeta\in C^1(Q)$ 成立。
\item 对每个非减 $h\in C(\mathbb{R})$，$\zeta\in C^1([0,T];C^1_c(\Omega))$，$\zeta\geqslant0$，及 $t\in(0,T]$，
\begin{align}
\begin{aligned}
&\iint_{Q_t}H(u_n)\partial_s\zeta\,\mathrm{d}x\,\mathrm{d}s
-\iint_{Q_t}h(u_n)\nabla\phi_n(u_n)\cdot\nabla\zeta\,\mathrm{d}x\,\mathrm{d}s\\
&\quad\geqslant\int_\Omega H(u_n(x,t))\zeta(x,t)\,\mathrm{d}x
-\int_\Omega H(u_{0n}(x))\zeta(x,0)\,\mathrm{d}x.
\end{aligned} \label{align:ch6-6.60}
\end{align}
若 $h(0)=0$，则\eqref{align:ch6-6.60}对所有 $\zeta\in C^1(Q)$，$\zeta\geqslant0$ 成立。
\end{enumerate}
\end{lemma}

\begin{proof}
因 $u_n$ 光滑，等式\eqref{align:ch6-6.59}由将 (Pn) 第一式乘以 $h(u_n)\zeta$（$h\in\operatorname{Lip}(\mathbb{R})$，$h(0)=0$，$\zeta\in C^1(Q)$）并在 $Q$ 上积分直接得到。此外，若 $\zeta\in C^1([0,T];C^1_c(\Omega))$，条件 $h(0)=0$ 可省略。故 (i) 得证。

因对所有 $n\in\mathbb{N}$ 有 $\nabla u_n\cdot\nabla\phi_n(u_n)\geqslant0$，由\eqref{align:ch6-6.59}对每个非减 $h\in\operatorname{Lip}(\mathbb{R})$ 得\eqref{align:ch6-6.60}。因每个非减 $h\in C(\mathbb{R})$ 可由一列非减 Lipschitz 函数局部逼近，结论随之成立。
\end{proof}

\begin{remark}\label{remark:ch06-6.5.2}
设 (A0)--(A1) 成立，$f\in C(\mathbb{R})$ 非减，$F_\phi$ 由\eqref{align:ch6-6.9}定义。在\eqref{align:ch6-6.60}中置 $g=f\circ\phi$，得不等式
\begin{align}
\begin{aligned}
&\int_\Omega F_\phi(u_n(x,t_2))\zeta(x,t_2)\,\mathrm{d}x
-\int_\Omega F_\phi(u_n(x,t_1))\zeta(x,t_1)\,\mathrm{d}x\\
&\quad\leqslant\int_{t_1}^{t_2}\int_\Omega\{F_\phi(u_n)\partial_t\zeta
-f(\phi(u_n))\nabla[\phi_n(u_n)]\cdot\nabla\zeta\}\,\mathrm{d}x\,\mathrm{d}t
\end{aligned} \label{align:ch6-6.61}
\end{align}
对所有 $t_1,t_2\in[0,T]$，$t_1<t_2$，及 $\zeta\in C^1([0,T];C^1_c(\Omega))$，$\zeta\geqslant0$ 成立。
\end{remark}

\begin{lemma}\label{lemma:ch06-6.5.5}
设 (A0)--(A1) 成立，$\{u_n\}$ 为问题 (Pn) 的解序列。则：
\begin{enumerate}[(i)]
\item 存在 $C_1=C_1(\lambda_N(\Omega),\|u_0\|_{\mathscr{R}_f(\Omega)})>0$，使对所有 $n\in\mathbb{N}$，
\begin{align}
\iint_Q\frac{|\nabla u_n|^2}{n^2}\,\mathrm{d}x\,\mathrm{d}t\leqslant\frac{C_1}{\sqrt{n}}; \label{align:ch6-6.62}
\end{align}
\item 存在 $C_2=C_2(\lambda_N(\Omega),\|\phi\|_\infty,\|u_0\|_{\mathscr{R}_f(\Omega)})>0$，使对所有 $n\in\mathbb{N}$，
\begin{align}
\iint_Q|\nabla[\phi_n(u_n)]|^2\,\mathrm{d}x\,\mathrm{d}t\leqslant C_2. \label{align:ch6-6.63}
\end{align}
\end{enumerate}
\end{lemma}

\begin{proof}
将 $\partial_t u_n=\Delta\phi_n(u_n)$ 乘以 $u_n/n$ 并在 $Q$ 上积分，直接得
\begin{align*}
\frac{1}{2n}\int_\Omega u_n^2(x,T)\,\mathrm{d}x-\frac{1}{2n}\int_\Omega u_{0n}^2(x)\,\mathrm{d}x
\leqslant-\iint_Q\frac{|\nabla u_n|^2}{n^2}\,\mathrm{d}x\,\mathrm{d}t.
\end{align*}
由此不等式及\eqref{align:ch6-6.49a}中第二个不等式，得
\begin{align*}
\iint_Q\frac{|\nabla u_n|^2}{n^2}\,\mathrm{d}x\,\mathrm{d}t
\leqslant\frac{1}{2n}\int_\Omega u_{0n}^2\,\mathrm{d}x
\leqslant\frac{M^2\lambda_N(\Omega)}{2\sqrt{n}}\|u_0\|^2_{\mathscr{R}_f(\Omega)}
=:\frac{C_1}{\sqrt{n}}.
\end{align*}
这证明了\eqref{align:ch6-6.62}。结合\eqref{align:ch6-6.57}与\eqref{align:ch6-6.62}直接得\eqref{align:ch6-6.63}，结论成立。
\end{proof}

\begin{lemma}\label{lemma:ch06-6.5.6}
设 (A1) 成立，$\{u_n\}$ 为问题 (Pn) 的解序列。设 $h\in\operatorname{Lip}(\mathbb{R})$ 有界且非减，$h(0)=0$，$H$ 为\eqref{align:ch6-6.58}中的函数。则对任意 $\rho\in C^1_c(\Omega)$，存在 $M_{h,\rho}>0$，使对所有 $n\in\mathbb{N}$，
\begin{align}
\int_0^T\left|\frac{\mathrm{d}}{\mathrm{d}t}\int_\Omega H(u_n(x,t))\rho\,\mathrm{d}x\right|\mathrm{d}t\leqslant M_{h,\rho}. \label{align:ch6-6.64}
\end{align}
\end{lemma}

\begin{proof}
对任意 $\rho\in C^1_c(\Omega)$，置
\begin{align*}
S_{n,\rho}(t):=\int_\Omega H(u_n(x,t))\rho(x)\,\mathrm{d}x \quad(t\in(0,T)).
\end{align*}
在\eqref{align:ch6-6.59}中取 $\zeta(x,t)=\beta(t)\rho(x)$，$\beta\in C^1_c(0,T)$，表明 $S_{n,\rho}$ 有弱导数
\begin{align}
\begin{aligned}
S'_{n,\rho}(t)
={}&-\int_\Omega h(u_n(x,t))\nabla\phi_n(u_n(x,t))\cdot\nabla\rho(x)\,\mathrm{d}x\\
&-\int_\Omega h'(u_n(x,t))\nabla u_n(x,t)\cdot\nabla\phi_n(u_n(x,t))\rho(x)\,\mathrm{d}x
\end{aligned} \label{align:ch6-6.65}
\end{align}
对几乎所有 $t\in(0,T)$。由\eqref{align:ch6-6.63}，
\begin{align}
\iint_Q|h(u_n)\nabla\phi_n(u_n)\cdot\nabla\rho|\,\mathrm{d}x\,\mathrm{d}t
\leqslant\|h\|_{L^\infty(\mathbb{R})}\|\rho\|_{C^1(\overline\Omega)}
\iint_Q|\nabla\phi_n(u_n)|\,\mathrm{d}x\,\mathrm{d}t
\leqslant C_2\|h\|_{L^\infty(\mathbb{R})}\|\rho\|_{C^1(\overline\Omega)}. \label{align:ch6-6.66}
\end{align}
估计\eqref{align:ch6-6.65}右端第二项。在\eqref{align:ch6-6.59}中取 $\zeta=1$，$t=T$，利用\eqref{align:ch6-6.49a}中第一个不等式，直接得
\begin{align}
\begin{aligned}
0\leqslant\iint_Q h'(u_n)\nabla u_n\cdot\nabla\phi_n(u_n)\,\mathrm{d}x\,\mathrm{d}t
\leqslant\int_\Omega H(u_{0n})\,\mathrm{d}x
\leqslant\|h\|_{L^\infty(\mathbb{R})}\|u_{0n}\|_{L^1(\Omega)}
\leqslant\|h\|_{L^\infty(\mathbb{R})}\|u_0\|_{\mathscr{R}_f(\Omega)}
\end{aligned} \label{align:ch6-6.67}
\end{align}
（注意 $H(u_n(\cdot,T))\geqslant0$ 几乎处处在 $\Omega$ 上，因 $h$ 非减且 $h(0)=0$）。由\eqref{align:ch6-6.65}--\eqref{align:ch6-6.67}有 $S'_{n,\rho}\in L^1(0,T)$，且
\begin{align}
\int_0^T|S'_{n,\rho}(t)|\,\mathrm{d}t
\leqslant(C_2\|\rho\|_{C^1(\overline\Omega)}+\|\rho\|_{C(\overline\Omega)}\|u_0\|_{\mathscr{R}_f(\Omega)})
\|h\|_{L^\infty(\mathbb{R})}T. \label{align:ch6-6.68}
\end{align}
结论随之成立。
\end{proof}

\begin{lemma}\label{lemma:ch06-6.5.7}
设 (A0)--(A1) 成立，$\{u_n\}$ 为问题 (Pn) 的解序列。则：
\begin{enumerate}[(i)]
\item 对每个 $t\in(0,T)$，
\begin{align}
\int_\Omega|\nabla[\phi_n(u_n)]|^2(x,s)\,\mathrm{d}x\leqslant\frac{C_2}{t}
\quad\text{对所有 }s\in(t,T]\text{ 及 }n\in\mathbb{N} \label{align:ch6-6.69}
\end{align}
其中 $C_2>0$ 由\cref{lemma:ch06-6.5.5}(ii) 给出；
\item 对所有 $t\in(0,T)$ 及 $n\in\mathbb{N}$，
\begin{align}
\int_t^T\int_\Omega[\phi'_n(u_n)](\partial_\tau u_n)^2\,\mathrm{d}x\,\mathrm{d}\tau
\leqslant\frac{C_2}{2t}. \label{align:ch6-6.70a}
\end{align}
进而，若 (A2) 成立，则对每个 $t\in(0,T)$，
\begin{align}
\int_t^T\int_\Omega(\partial_\tau[\phi_n(u_n)])^2\,\mathrm{d}x\,\mathrm{d}\tau
\leqslant\frac{C_2(1+\|\phi'\|_\infty)}{2t}. \label{align:ch6-6.70b}
\end{align}
\end{enumerate}
\end{lemma}

\begin{proof}
(i) 对任意固定 $t\in(0,T)$，置 $h(s):=\frac{s}{t}\chi_{[0,t)}(s)+\chi_{[t,T)}(s)$（$s\in(t,T]$）。将 $\partial_t u_n=\Delta\phi_n(u_n)$ 乘以 $h(t)\partial_t[\phi_n(u_n)]$ 并在 $Q_s=\Omega\times(0,s)$ 上积分，对每个 $s\in(t,T)$，得
\begin{align}
\begin{aligned}
0&\leqslant\iint_{Q_s}h(\tau)\phi'_n(u_n)(\partial_\tau u_n)^2\,\mathrm{d}x\,\mathrm{d}\tau\\
&=-\int_0^s\mathrm{d}\tau\,h(\tau)\int_\Omega\partial_\tau\{\nabla[\phi_n(u_n)]\}\cdot\nabla[\phi_n(u_n)]\,\mathrm{d}x\\
&=-\frac{1}{2}\int_0^s h(\tau)\frac{\mathrm{d}}{\mathrm{d}\tau}\left(\int_\Omega|\nabla[\phi_n(u_n)]|^2\,\mathrm{d}x\right)\mathrm{d}\tau\\
&=-\frac{1}{2}\int_\Omega|\nabla[\phi_n(u_n)]|^2(x,s)\,\mathrm{d}x
+\frac{1}{2t}\int_0^t\mathrm{d}\tau\int_\Omega|\nabla[\phi_n(u_n)]|^2\,\mathrm{d}x.
\end{aligned} \label{align:ch6-6.71}
\end{align}
利用\eqref{align:ch6-6.63}，由上述不等式得\eqref{align:ch6-6.69}。故 (i) 得证。

(ii) 在\eqref{align:ch6-6.71}中取 $s=T$，得
\begin{align*}
\int_t^T\int_\Omega[\phi'_n(u_n)](\partial_\tau u_n)^2\,\mathrm{d}x\,\mathrm{d}\tau
\leqslant\iint_Q h(\tau)\phi'_n(u_n)(\partial_\tau u_n)^2\,\mathrm{d}x\,\mathrm{d}\tau
\leqslant\frac{1}{2t}\int_0^t\mathrm{d}\tau\int_\Omega|\nabla[\phi_n(u_n)]|^2\,\mathrm{d}x
\leqslant\frac{C_2}{2t}.
\end{align*}
这证明了\eqref{align:ch6-6.70a}，\eqref{align:ch6-6.70b}直接成立。证明完成。
\end{proof}

\begin{remark}\label{remark:ch06-6.5.3}
因
\begin{align*}
|\nabla[\phi_n(u_n)]|^2=\left(\phi'(u_n)+\frac{1}{n}\right)^2|\nabla u_n|^2
\geqslant(\phi'(u_n))^2|\nabla u_n|^2=|\nabla\phi(u_n)|^2,
\end{align*}
由\eqref{align:ch6-6.69}，对所有 $0<t<s\leqslant T$，得
\begin{align}
\int_\Omega|\nabla\phi(u_n)|^2(x,s)\,\mathrm{d}x\leqslant\frac{C_2}{t}. \label{align:ch6-6.72}
\end{align}
类似地，因 $|\partial_\tau\phi(u_n)|^2\leqslant|\partial_\tau[\phi_n(u_n)]|^2$，由\eqref{align:ch6-6.70b}得
\begin{align}
\int_t^T\int_\Omega(\partial_\tau[\phi(u_n)])^2\,\mathrm{d}x\,\mathrm{d}\tau
\leqslant\frac{C_2(1+\|\phi'\|_\infty)}{2t}. \label{align:ch6-6.73}
\end{align}
\end{remark}

\begin{lemma}\label{lemma:ch06-6.5.8}
设 (A0)--(A1) 成立，$\{u_n\}$ 为问题 (Pn) 的解序列。则：
\begin{enumerate}[(i)]
\item 对所有 $k>0$ 及 $n\in\mathbb{N}$，有
\begin{align}
\iint_{\{\pm u_n>k\}}|\nabla\phi(u_n)|^2\,\mathrm{d}x\,\mathrm{d}t
\leqslant\|u^\pm_0\|_{\mathscr{R}_f(\Omega)}[\pm\phi_{\pm\infty}\mp\phi(\pm k)], \label{align:ch6-6.74}
\end{align}
\item 存在 $M_0>0$，使对所有 $k>0$ 及 $n\in\mathbb{N}$，
\begin{align}
\iint_{\{\pm u_n>k\}}|\nabla\phi_n(u_n)|^2\,\mathrm{d}x\,\mathrm{d}t
\leqslant M_0[\pm\phi_{\pm\infty}\mp\phi(\pm k)+n^{-1/2}]. \label{align:ch6-6.75}
\end{align}
\end{enumerate}
\end{lemma}

\begin{proof}
我们仅证\eqref{align:ch6-6.74}中上号情形，另一情形类似。对任意 $k>0$，置
\begin{align*}
F(y):=\chi_{(k,+\infty)}(y)\int_k^y[\phi(z)-\phi(k)]\,\mathrm{d}z \quad(y\in\mathbb{R}).
\end{align*}
于是 $F'\in W^{1,\infty}_{\mathrm{loc}}(\mathbb{R})$，且
\begin{align}
0\leqslant F(y)\leqslant[\phi_{+\infty}-\phi(k)](y-k)\chi_{(k,+\infty)}(y). \label{align:ch6-6.76}
\end{align}
将 $\partial_t u_n=\Delta\phi_n(u_n)$ 乘以 $F'(u_n)$ 并在 $Q$ 上积分，由\eqref{align:ch6-6.76}与\eqref{align:ch6-6.49a}得
\begin{align*}
\begin{aligned}
&\iint_{\{u_n>k\}}|\nabla\phi(u_n)|^2\,\mathrm{d}x\,\mathrm{d}t
+\frac{1}{n}\iint_{\{u_n>k\}}\phi'(u_n)|\nabla u_n|^2\,\mathrm{d}x\,\mathrm{d}t\\
&\quad=\int_\Omega F(u_{0n}(x))\,\mathrm{d}x-\int_\Omega F(u_n(x,T))\,\mathrm{d}x\\
&\quad\leqslant[\phi_{+\infty}-\phi(k)]\int_{\{u_{0n}\geqslant k\}}[u_{0n}(x)-k]\,\mathrm{d}x
\leqslant\|u^+_0\|_{\mathscr{R}_f(\Omega)}[\phi_{+\infty}-\phi(k)].
\end{aligned}
\end{align*}
于是\eqref{align:ch6-6.74}得证。关于\eqref{align:ch6-6.75}，由\eqref{align:ch6-6.74}与\eqref{align:ch6-6.62}有
\begin{align*}
\begin{aligned}
\iint_{\{\pm u_n>k\}}|\nabla\phi_n(u_n)|^2\,\mathrm{d}x\,\mathrm{d}t
&\leqslant2\iint_{\{\pm u_n>k\}}\left(|\nabla\phi(u_n)|^2+\frac{|\nabla u_n|^2}{n^2}\right)\mathrm{d}x\,\mathrm{d}t\\
&\leqslant2\|u_0\|_{\mathscr{R}_f(\Omega)}[\pm\phi_{\pm\infty}\mp\phi(\pm k)]
+2C_1n^{-1/2}
\end{aligned}
\end{align*}
其中 $C_1>0$ 如\cref{lemma:ch06-6.5.5}(i)。置 $M_0:=2\max\{\|u_0\|_{\mathscr{R}_f(\Omega)},C_1\}$，得\eqref{align:ch6-6.75}，结论成立。
\end{proof}

\section{存在性结果的证明：令 $n\to\infty$}\label{section:ch06-令n趋于无穷}

\cref{theorem:ch06-6.4.1}的证明依赖如下结果，其证明需要若干步骤。

\begin{theorem}\label{theorem:ch06-6.6.1}
设 (A0)--(A2) 成立，$\{u_n\}$ 为问题 (Pn) 的解序列。则存在子列 $\{u_{n_k}\}\subseteq\{u_n\}$，具有如下性质：
\begin{enumerate}[(i)]
\item 存在 $u\in L^\infty_{w^*}(0,T;\mathscr{R}_f(\Omega))$，使 $\phi(u_r)\in L^2(0,T;H^1_0(\Omega))$，且
\begin{align}
u_{n_k}\stackrel{*}{\rightharpoonup}u \quad\text{在 }L^\infty_{w^*}(0,T;\mathscr{R}_f(\Omega)); \label{align:ch6-6.77}
\end{align}
\item 与 $\{u_{n_k}\}$ 相伴的 Young 测度序列 $\{\nu_{n_k}\}\subseteq\mathcal{Y}^+(Q;\mathbb{R})$ 窄收敛于 Young 测度 $\nu\in\mathcal{Y}^+(Q;\mathbb{R})$，且对几乎所有 $(x,t)\in Q$，有
\begin{align}
u_r(x,t)=\int_\mathbb{R}y\,\mathrm{d}\nu_{(x,t)}(y), \label{align:ch6-6.78}
\end{align}
\begin{align}
\operatorname{supp}\nu_{(x,t)}\subseteq\phi^{-1}(\{\phi(u_r(x,t))\}), \label{align:ch6-6.79}
\end{align}
其中 $\{\nu_{(x,t)}\}_{(x,t)\in Q}$ 为 $\nu$ 的离散分解；
\item 对每个满足\cref{proposition:ch05-函数弱星}的 $f\in C(\mathbb{R})$ 及所有 $\zeta\in C([0,T];C_c(\Omega))$，有
\begin{align}
\begin{aligned}
\lim_{k\to\infty}\iint_Q f(u_{n_k})\zeta\,\mathrm{d}x\,\mathrm{d}t
={}&\iint_Q f^*\zeta\,\mathrm{d}x\,\mathrm{d}t
+M^+_f\int_0^T\langle u^+_s(\cdot,t),\zeta(\cdot,t)\rangle\,\mathrm{d}t\\
&-M^-_f\int_0^T\langle u^-_s(\cdot,t),\zeta(\cdot,t)\rangle\,\mathrm{d}t,
\end{aligned} \label{align:ch6-6.80}
\end{align}
其中 $f^*\in L^\infty(0,T;L^1(\Omega))$ 定义为
\begin{align}
f^*(x,t):=\int_\mathbb{R}f(y)\,\mathrm{d}\nu_{(x,t)}(y)
\quad\text{对几乎所有 }(x,t)\in Q. \label{align:ch6-6.81}
\end{align}
\end{enumerate}
\end{theorem}

为证\cref{theorem:ch06-6.6.1}，需要如下结果。

\begin{proposition}\label{proposition:ch06-6.6.2}
设 (A0)--(A2) 成立，$\{u_n\}$ 为问题 (Pn) 的解序列。则存在子列 $\{u_{n_k}\}\subseteq\{u_n\}$ 及 Young 测度 $\nu\in\mathcal{Y}^+(Q;\mathbb{R})$，使：
\begin{align}
\phi(u_{n_k})\stackrel{*}{\rightharpoonup}\phi^* \text{ 在 }L^\infty(Q),\quad
\phi(u_{n_k})\rightharpoonup\phi^* \text{ 在 }L^2(Q), \label{align:ch6-6.82}
\end{align}
其中
\begin{align}
\phi^*(x,t):=\int_\mathbb{R}\phi(y)\,\mathrm{d}\nu_{(x,t)}(y)\in L^2(0,T;H^1_0(\Omega)), \label{align:ch6-6.83}
\end{align}
$\{\nu_{(x,t)}\}_{(x,t)\in Q}$ 为 $\nu$ 的离散分解；
\begin{align}
\nabla\phi(u_{n_k}),\ \nabla\phi_{n_k}(u_{n_k})\rightharpoonup\nabla\phi^* \quad\text{在 }[L^2(Q)]^N; \label{align:ch6-6.84}
\end{align}
对任意 $\tau\in(0,T)$，
\begin{align}
\partial_t[\phi(u_{n_k})],\ \partial_t[\phi_{n_k}(u_{n_k})]\rightharpoonup\partial_t\phi^*
\quad\text{在 }L^2(\Omega\times(\tau,T)), \label{align:ch6-6.85}
\end{align}
\begin{align}
\phi(u_{n_k}),\ \phi_{n_k}(u_{n_k})\rightharpoonup\phi^* \quad\text{在 }H^1(\Omega\times(\tau,T)), \label{align:ch6-6.86}
\end{align}
\begin{align}
\phi(u_{n_k}),\ \phi_{n_k}(u_{n_k})\to\phi^* \text{ 在 }Q\text{ 中几乎处处}; \label{align:ch6-6.87}
\end{align}
且对几乎所有 $t\in(0,T)$，
\begin{align}
\phi(u_{n_k}(\cdot,t)),\ \phi_{n_k}(u_{n_k}(\cdot,t))\to\phi^*(\cdot,t)
\quad\text{在 }\Omega\text{ 中几乎处处}, \label{align:ch6-6.88}
\end{align}
\begin{align}
\phi(u_{n_k}(\cdot,t))\rightharpoonup\phi^*(\cdot,t) \quad\text{在 }H^1_0(\Omega). \label{align:ch6-6.89}
\end{align}
\end{proposition}

\begin{proof}
因序列 $\{\phi(u_n)\}$ 在 $L^\infty(Q)$ 中有界，由 Banach--Alaoglu 定理存在子列 $\{u_{n_k}\}\subseteq\{u_n\}$ 及 $v\in L^\infty(Q)$，使 $\phi(u_{n_k})\stackrel{*}{\rightharpoonup}v$ 在 $L^\infty(Q)$。同时序列 $\{u_n\}$ 在 $L^1(Q)$ 中有界，故其相伴 Young 测度序列 $\{\nu_n\}$ 在 $Q\times\mathbb{R}$ 上 Young 测度的窄拓扑中相对紧，极限 Young 测度 $\nu$ 的存在性直接得到。又因 $Q$ 有界，子列 $\{\phi(u_{n_k})\}$ 在 $L^1(Q)$ 中有界且一致可积（见\cref{lemma:ch02-有界一致可积}）。故由\cref{proposition:ch05-弱紧YoungL1}存在 $\{\phi(u_{n_k})\}$ 的子列（不重新标号），使 $\phi(u_{n_k})\rightharpoonup\phi^*$ 在 $L^1(Q)$，$\phi^*$ 如\eqref{align:ch6-6.83}。于是 $v=\phi^*$，\eqref{align:ch6-6.82}得证。

关于\eqref{align:ch6-6.84}，注意由\eqref{align:ch6-6.57}与\eqref{align:ch6-6.63}，序列 $\{\nabla\phi(u_{n_k})\}$ 与 $\{\nabla\phi_{n_k}(u_{n_k})\}$ 在 $[L^2(Q)]^N$ 中弱收敛。由\eqref{align:ch6-6.82}得 $\nabla\phi(u_{n_k})\rightharpoonup\nabla\phi^*$ 在 $[L^2(Q)]^N$。因
\begin{align}
|\phi_{n_k}(u_{n_k})-\phi(u_{n_k})|=\frac{|u_{n_k}|}{n_k}\to0 \text{ 在 }L^1(Q) \label{align:ch6-6.90}
\end{align}
（见\eqref{align:ch6-6.56}），也得\eqref{align:ch6-6.84}中第二个收敛。由\eqref{align:ch6-6.84}、\eqref{align:ch6-6.82}中第二个收敛及范数半连续性，有 $\phi^*\in L^2(0,T;H^1_0(\Omega))$。

现在证\eqref{align:ch6-6.85}。由\eqref{align:ch6-6.82}中第二个收敛，对所有 $\zeta\in C^1_c(Q)$，
\begin{align}
\lim_{k\to\infty}\iint_Q\partial_t[\phi(u_{n_k})]\zeta\,\mathrm{d}x\,\mathrm{d}t
=-\lim_{k\to\infty}\iint_Q\phi(u_{n_k})\partial_t\zeta\,\mathrm{d}x\,\mathrm{d}t
=-\iint_Q\phi^*\partial_t\zeta\,\mathrm{d}x\,\mathrm{d}t. \label{align:ch6-6.91}
\end{align}
另一方面，由\eqref{align:ch6-6.70b}，对任意 $\tau\in(0,T)$，序列 $\{\partial_t[\phi(u_{n_k})]\}$ 在 $L^2(\Omega\times(\tau,T))$ 中有界，故存在 $\{\partial_t[\phi(u_{n_k})]\}$ 的子列（不重新标号）及 $g_\tau\in L^2(\Omega\times(\tau,T))$，使对所有 $\zeta\in L^2(\Omega\times(\tau,T))$，
\begin{align}
\lim_{k\to\infty}\iint_{\Omega\times(\tau,T)}\partial_t[\phi(u_{n_k})]\zeta\,\mathrm{d}x\,\mathrm{d}t
=\iint_{\Omega\times(\tau,T)}g_\tau\zeta\,\mathrm{d}x\,\mathrm{d}t. \label{align:ch6-6.92}
\end{align}
由\eqref{align:ch6-6.91}--\eqref{align:ch6-6.92}，对每个 $\zeta\in C^1_c(\Omega\times(\tau,T))$，有
\begin{align}
\iint_{\Omega\times(\tau,T)}g_\tau\zeta\,\mathrm{d}x\,\mathrm{d}t
=-\iint_Q\phi^*\partial_t\zeta\,\mathrm{d}x\,\mathrm{d}t. \label{align:ch6-6.93}
\end{align}
于是分布导数 $\partial_t\phi^*$ 可等同于 $L^2(\Omega\times(\tau,T))$ 中一元，故\eqref{align:ch6-6.85}中第一个收敛成立。第二个类似证明。

\eqref{align:ch6-6.86}中的收敛由\eqref{align:ch6-6.84}--\eqref{align:ch6-6.85}显然成立。故序列 $\{\phi(u_{n_k})\}$ 与 $\{\phi_{n_k}(u_{n_k})\}$ 均在 $H^1(\Omega\times(\tau,T))$ 中有界。由嵌入定理（可能抽取子列，为简便不重新标号），对每个 $\tau\in(0,T)$，有 $\phi_{n_k}(u_{n_k}),\phi(u_{n_k})\to\phi^*$ 在 $\Omega\times(\tau,T)$ 中几乎处处。由 $\tau$ 的任意性及标准对角论证，可构造子列（为简便不重新标号），使\eqref{align:ch6-6.87}中两个收敛同时成立。

最后，\eqref{align:ch6-6.88}中的收敛由\eqref{align:ch6-6.87}得到，而\eqref{align:ch6-6.89}是\eqref{align:ch6-6.88}与\eqref{align:ch6-6.69}的直接推论。证明完成。
\end{proof}

\begin{proposition}\label{proposition:ch06-6.6.3}
设 (A0)--(A1) 成立，$\{u_n\}$ 为问题 (Pn) 的解序列，$\nu\in\mathcal{Y}^+(Q;\mathbb{R})$ 为\cref{proposition:ch06-6.6.2}中给出的 Young 测度。则存在子列 $\{u_{n_k}\}\subseteq\{u_n\}$ 及 $\sigma^\pm\in L^\infty_{w^*}(0,T;\mathscr{R}^+_f(\Omega))$，使
\begin{align}
u^\pm_{n_k}\stackrel{*}{\rightharpoonup}u_{b,\pm}+\sigma^\pm
\quad\text{在 }L^\infty_{w^*}(0,T;\mathscr{R}_f(\Omega)), \label{align:ch6-6.94}
\end{align}
其中 $u_{b,\pm}\in L^\infty(0,T;L^1(\Omega))$，
\begin{align}
u_{b,\pm}(x,t):=\int_\mathbb{R}y^\pm\,\mathrm{d}\nu_{(x,t)}(y)\geqslant0
\quad\text{对几乎所有 }(x,t)\in Q. \label{align:ch6-6.95}
\end{align}
\end{proposition}

\begin{proof}
由不等式\eqref{align:ch6-6.54}，序列 $\{u_n\}$ 在 $L^\infty(0,T;L^1(\Omega))\subseteq L^\infty_{w^*}(0,T;\mathscr{R}_f(\Omega))$ 中有界。于是由\cref{theorem:ch05-弱星紧性}与 Banach--Alaoglu 定理，存在 $\mu^\pm\in L^\infty_{w^*}(0,T;\mathscr{R}^+_f(\Omega))$，使
\begin{align}
u^\pm_{n_k}\stackrel{*}{\rightharpoonup}\mu^\pm
\quad\text{在 }L^\infty_{w^*}(0,T;\mathscr{R}_f(\Omega)). \label{align:ch6-6.96}
\end{align}
另一方面，因序列 $\{u_{n_k}\}$ 在 $L^1(Q)$ 中有界，由\cref{remark:ch05-biting注}存在 $\{u_{n_k}\}$ 的子列（不重新标号）及 $\sigma^\pm\in\mathscr{R}^+_f(Q)$，使
\begin{align}
u^\pm_{n_k}\stackrel{*}{\rightharpoonup}u_{b,\pm}+\sigma^\pm
\quad\text{在 }\mathscr{R}_f(Q). \label{align:ch6-6.97}
\end{align}
由\eqref{align:ch6-6.96}--\eqref{align:ch6-6.97}，$\mu^\pm=u_{b,\pm}+\sigma^\pm$，结论成立。
\end{proof}

\begin{proposition}\label{proposition:ch06-6.6.4}
设 (A0)--(A1) 成立，$\{u_n\}$ 为问题 (Pn) 的解序列。设 $\{u_{n_k}\}\subseteq\{u_n\}$，$\nu\in\mathcal{Y}^+(Q;\mathbb{R})$，及 $\sigma^\pm\in L^\infty_{w^*}(0,T;\mathscr{R}^+_f(\Omega))$ 由\cref{proposition:ch06-6.6.3}给出。则对每个满足\cref{proposition:ch05-函数弱星}的 $f\in C(\mathbb{R})$ 及所有 $\zeta\in C([0,T];C_c(\Omega))$，有
\begin{align}
\begin{aligned}
\lim_{k\to\infty}\iint_Q f(u_{n_k})\zeta\,\mathrm{d}x\,\mathrm{d}t
={}&\iint_Q f^*\zeta\,\mathrm{d}x\,\mathrm{d}t
+M^+_f\int_0^T\langle\sigma^+(\cdot,t),\zeta(\cdot,t)\rangle\,\mathrm{d}t\\
&-M^-_f\int_0^T\langle\sigma^-(\cdot,t),\zeta(\cdot,t)\rangle\,\mathrm{d}t
\end{aligned} \label{align:ch6-6.98}
\end{align}
其中 $f^*\in L^\infty(0,T;L^1(\Omega))$ 定义为
\begin{align}
f^*(x,t):=\int_\mathbb{R}f(y)\,\mathrm{d}\nu_{(x,t)}(y)
\quad\text{对几乎所有 }(x,t)\in Q. \label{align:ch6-6.99}
\end{align}
\end{proposition}

\begin{proof}
由\cref{proposition:ch05-函数弱星}，存在 $L>0$，使对所有 $y\in\mathbb{R}$，$|f(y)|\leqslant L(1+|y|)$，故由不等式\eqref{align:ch6-6.54}，
\begin{align*}
\sup_{n\in\mathbb{N}}\|f(u_n)\|_{L^\infty(0,T;L^1(\Omega))}
=\sup_{n\in\mathbb{N}}\|f(u_n)\|_{L^\infty_{w^*}(0,T;\mathscr{R}_f(\Omega))}
\leqslant L[\lambda_N(\Omega)+\|u_0\|_{\mathscr{R}_f(\Omega)}].
\end{align*}
由 Banach--Alaoglu 定理，存在子列 $\{f(u_{n_k})\}\subseteq\{f(u_n)\}$ 及 $\mu\in L^\infty_{w^*}(0,T;\mathscr{R}_f(\Omega))$，使
\begin{align}
f(u_{n_k})\stackrel{*}{\rightharpoonup}\mu
\quad\text{在 }L^\infty_{w^*}(0,T;\mathscr{R}_f(\Omega)). \label{align:ch6-6.100}
\end{align}
此外，如\cref{proposition:ch05-集中测度极限}证明中那样论证，得
\begin{align}
f(u_{n_k})\stackrel{*}{\rightharpoonup}f^*+M^+_f\sigma^+-M^-_f\sigma^-
\quad\text{在 }\mathscr{R}_f(Q). \label{align:ch6-6.101}
\end{align}
由\eqref{align:ch6-6.100}--\eqref{align:ch6-6.101}，结论成立。
\end{proof}

置
\begin{align}
u_b(x,t):=u_{b,+}(x,t)-u_{b,-}(x,t)=\int_\mathbb{R}y\,\mathrm{d}\nu_{(x,t)}(y)
\quad\text{对几乎所有 }(x,t)\in Q, \label{align:ch6-6.102a}
\end{align}
\begin{align}
u:=u_b+\sigma^+-\sigma^-. \label{align:ch6-6.102b}
\end{align}

由\cref{proposition:ch06-6.6.3}与\cref{proposition:ch06-6.6.4}，有如下结果。

\begin{proposition}\label{proposition:ch06-6.6.5}
设 (A0)--(A1) 成立，$\{u_n\}$ 为问题 (Pn) 的解序列。设 $\{u_{n_k}\}\subseteq\{u_n\}$ 及 $\sigma^\pm\in L^\infty_{w^*}(0,T;\mathscr{R}^+_f(\Omega))$ 由\cref{proposition:ch06-6.6.3}给出，$u\in L^\infty_{w^*}(0,T;\mathscr{R}_f(\Omega))$ 由\eqref{align:ch6-6.102b}定义。则\eqref{align:ch6-6.77}中的收敛成立。进而，
\begin{align}
u_r=u_b \quad\text{在 }Q\text{ 中几乎处处}, \label{align:ch6-6.103a}
\end{align}
\begin{align}
u^\pm_s=\sigma^\pm \quad\text{在 }L^\infty_{w^*}(0,T;\mathscr{R}_f(\Omega)). \label{align:ch6-6.103b}
\end{align}
\end{proposition}

\begin{proof}
\eqref{align:ch6-6.77}中的收敛是\eqref{align:ch6-6.94}--\eqref{align:ch6-6.95}与\eqref{align:ch6-6.102a}的直接推论。此外，其余断言\eqref{align:ch6-6.103a}--\eqref{align:ch6-6.103b}将由证明 $\sigma^\pm$ 相互奇异且关于 Lebesgue 测度奇异而得到。

为此，证明如下断言。

\begin{quote}
对任意 $k\in\mathbb{N}$，$p>0$，及 $\zeta\in C^1([0,T];C^2_c(\Omega))$，$\zeta\geqslant0$ 在 $Q$ 中，$\zeta(\cdot,T)=0$，有
\begin{align}
\begin{aligned}
-\iint_Q[u_{n_k}\mp p]^\pm\partial_t\zeta\,\mathrm{d}x\,\mathrm{d}t
\leqslant{}&\int_\Omega[u_{0n_k}\mp p]^\pm\zeta(x,0)\,\mathrm{d}x\\
&+\|\zeta\|_{C^1(\overline\Omega)}\left(\frac{C_2T\|u_0\|_{\mathscr{R}_f(\Omega)}}{p}\right)^{1/2},
\end{aligned} \label{align:ch6-6.104}
\end{align}
其中 $C_2=C_2(\lambda_N(\Omega),\|\phi\|_\infty,\|u_0\|_{\mathscr{R}_f(\Omega)})>0$ 为\cref{lemma:ch06-6.5.5}(ii) 中给出的常数。
\end{quote}

依赖上述断言，可证
\begin{align}
\sigma^\pm(\cdot,t)\leqslant u^\pm_{0s} \quad\text{在 }\mathscr{R}^+_f(\Omega)
\quad\text{对几乎所有 }t\in(0,T). \label{align:ch6-6.105}
\end{align}
因测度 $\sigma^\pm$ 非负，这蕴含 $\sigma^\pm(\cdot,t)$ 关于 $u^\pm_{0s}$ 绝对连续，对几乎所有 $t\in(0,T)$。结论随之成立。

我们仅证\eqref{align:ch6-6.105}中 $+$ 号情形，$-$ 号证明类似。在\eqref{align:ch6-6.51}中取 $f(y)=[y-p]_+$（$y\in\mathbb{R}$，$p>0$），得
\begin{align}
[u_{0n_k}-p]^+\stackrel{*}{\rightharpoonup}[u_{0r}-p]^++u^+_{0s}
\quad\text{在 }\mathscr{R}_f(\Omega), \label{align:ch6-6.106}
\end{align}
而由\eqref{align:ch6-6.98}--\eqref{align:ch6-6.99}中同样的 $f$ 选择，得
\begin{align*}
\begin{aligned}
\lim_{k\to\infty}\iint_Q[u_{n_k}-p]^+\partial_t\zeta\,\mathrm{d}x\,\mathrm{d}t
={}&\iint_Q\left(\int_\mathbb{R}[y-p]^+\,\mathrm{d}\nu_{(x,t)}(y)\right)\partial_t\zeta\,\mathrm{d}x\,\mathrm{d}t\\
&+\int_0^T\langle\sigma^+(\cdot,t),\partial_t\zeta(\cdot,t)\rangle\,\mathrm{d}t
\end{aligned}
\end{align*}
对所有 $\zeta\in C^1([0,T];C^2_c(\Omega))$ 成立。由此等式，在\eqref{align:ch6-6.104}中取上号，令 $k\to\infty$，并利用\eqref{align:ch6-6.106}，得
\begin{align}
\begin{aligned}
&-\iint_Q\left(\int_\mathbb{R}[y-p]^+\,\mathrm{d}\nu_{(x,t)}(y)\right)\partial_t\zeta\,\mathrm{d}x\,\mathrm{d}t
-\int_0^T\langle\sigma^+(\cdot,t),\partial_t\zeta(\cdot,t)\rangle\,\mathrm{d}t\\
&\quad\leqslant\int_\Omega[u_{0r}-p]^+\zeta(x,0)\,\mathrm{d}x
+\langle u^+_{0s},\zeta(\cdot,0)\rangle
+\|\zeta\|_{C^1(\overline\Omega)}\left(\frac{C_2T\|u_0\|_{\mathscr{R}_f(\Omega)}}{p}\right)^{1/2}
\end{aligned} \label{align:ch6-6.107}
\end{align}
对每个 $\zeta\in C^1([0,T];C^2_c(\Omega))$，$\zeta\geqslant0$ 在 $Q$ 中，$\zeta(\cdot,T)=0$ 在 $\Omega$ 上成立。

因对每个 $y\in\mathbb{R}$ 有 $\lim_{p\to+\infty}[y-p]^+=0$，且对几乎所有 $(x,t)\in Q$，$[y-p]^+\leqslant y^+\in L^1(\mathbb{R},\mathscr{B},\nu_{(x,t)})$（见\cref{proposition:ch05-矩估计}(i)），由控制收敛定理，
\begin{align*}
\lim_{p\to\infty}\int_\mathbb{R}[y-p]^+\,\mathrm{d}\nu_{(x,t)}(y)=0
\quad\text{对几乎所有 }(x,t)\in Q.
\end{align*}
此外，对几乎所有 $(x,t)\in Q$，有 $\int_\mathbb{R}[y-p]^+\,\mathrm{d}\nu_{(x,t)}(y)\leqslant\int_\mathbb{R}y^+\,\mathrm{d}\nu_{(x,t)}(y)=u_{b,+}(x,t)$（见\eqref{align:ch6-6.95}）。因 $u_{b,+}\in L^\infty(0,T;L^1(\Omega))\subseteq L^1(Q)$，由控制收敛定理得
\begin{align}
\lim_{p\to\infty}\iint_Q\left(\int_\mathbb{R}[y-p]^+\,\mathrm{d}\nu_{(x,t)}(y)\right)\partial_t\zeta\,\mathrm{d}x\,\mathrm{d}t=0. \label{align:ch6-6.108a}
\end{align}
类似地，
\begin{align}
\lim_{p\to\infty}\int_\Omega[u_{0r}-p]^+\zeta(x,0)\,\mathrm{d}x=0. \label{align:ch6-6.108b}
\end{align}
由\eqref{align:ch6-6.108a}，如前在\eqref{align:ch6-6.107}中令 $p\to\infty$，对每个 $\zeta$ 得
\begin{align}
-\int_0^T\langle\sigma^+(\cdot,t),\partial_t\zeta(\cdot,t)\rangle\,\mathrm{d}t
\leqslant\langle u^+_{0s},\zeta(\cdot,0)\rangle. \label{align:ch6-6.109}
\end{align}

设 $t_0\in(0,T)$，对任意充分大的 $q>0$，置
\begin{align*}
h_q(t):=\chi_{[0,t_0-\frac{1}{q}]}(t)+q(t_0-t)\chi_{(t_0-\frac{1}{q},t_0]}(t) \quad(t\in(0,T)).
\end{align*}
由标准逼近论证，可在\eqref{align:ch6-6.109}中取 $\zeta(x,t)=\rho(x)h_q(t)$，$\rho\in C^2_c(\Omega)$，$\rho\geqslant0$。于是得
\begin{align*}
q\int_{t_0-\frac{1}{q}}^{t_0}\langle\sigma^+(\cdot,t),\rho\rangle\,\mathrm{d}t\leqslant\langle u^+_{0s},\rho\rangle,
\end{align*}
从而当 $q\to\infty$ 时，
\begin{align*}
\langle\sigma^+(\cdot,t_0),\rho\rangle\leqslant\langle u^+_{0s},\rho\rangle
\quad\text{对几乎所有 }t_0\in(0,T).
\end{align*}
由 $\rho$ 的任意性，\eqref{align:ch6-6.105}中 $+$ 号情形得证。

最后证明该断言。由标准逼近论证，对每个 $p>0$，可在\eqref{align:ch6-6.58}中取 $h=\chi_{(p,\infty)}$，从而在\eqref{align:ch6-6.59}中 $H(y)=[y-p]_+$。于是对每个非负 $\zeta\in C^1([0,T];C^1_c(\Omega))$，$\zeta(\cdot,T)=0$ 在 $\Omega$ 上，得（见\eqref{align:ch6-6.60}）
\begin{align}
-\iint_Q[u_{n_k}-p]^+\partial_t\zeta\,\mathrm{d}x\,\mathrm{d}t
\leqslant\int_\Omega[u_{0n_k}-p]^+\zeta(x,0)\,\mathrm{d}x
-\iint_{\{u_{n_k}>p\}}\nabla\phi_{n_k}(u_{n_k})\cdot\nabla\zeta\,\mathrm{d}x\,\mathrm{d}t. \label{align:ch6-6.110}
\end{align}
由\cref{lemma:ch06-6.5.5}(ii)，存在 $C_2=C_2(\lambda_N(\Omega),\|\phi\|_\infty,\|u_0\|_{\mathscr{R}_f(\Omega)})>0$，使
\begin{align*}
\iint_{\{u_{n_k}>p\}}|\nabla\phi_{n_k}(u_{n_k})||\nabla\zeta|\,\mathrm{d}x\,\mathrm{d}t
\leqslant C_2^{1/2}\|\zeta\|_{C^1(Q)}[\lambda_{N+1}(\{u_{n_k}>p\})]^{1/2},
\end{align*}
从而由 Chebyshev 不等式，
\begin{align}
\iint_{\{u_{n_k}>p\}}|\nabla\phi_{n_k}(u_{n_k})||\nabla\zeta|\,\mathrm{d}x\,\mathrm{d}t
\leqslant\|\zeta\|_{C^1(\overline\Omega)}\left(\frac{C_2\|u_{n_k}\|_{L^1(Q)}}{p}\right)^{1/2}. \label{align:ch6-6.111}
\end{align}
由\eqref{align:ch6-6.54}与\eqref{align:ch6-6.110}--\eqref{align:ch6-6.111}，得上号情形的\eqref{align:ch6-6.104}。结论随之成立。
\end{proof}

\begin{proposition}\label{proposition:ch06-6.6.6}
设 (A0)--(A2) 成立，$\{u_n\}$ 为问题 (Pn) 的解序列。设 $\{u_{n_k}\}\subseteq\{u_n\}$ 与 $\nu\in\mathcal{Y}^+(Q;\mathbb{R})$ 由\cref{proposition:ch06-6.6.3}给出，$u\in L^\infty_{w^*}(0,T;\mathscr{R}_f(\Omega))$ 由\eqref{align:ch6-6.102b}定义。则对几乎所有 $(x,t)\in Q$：
\begin{enumerate}[(i)]
\item 有
\begin{align}
\operatorname{supp}\nu_{(x,t)}\subseteq\phi^{-1}(\{\phi^*(x,t)\}) \label{align:ch6-6.112}
\end{align}
其中 $\phi^*$ 由\eqref{align:ch6-6.83}给出；
\item 有
\begin{align}
\phi^*(x,t)=\phi(u_r(x,t)),\quad
\phi(u_r)\in L^2(0,T;H^1_0(\Omega)), \label{align:ch6-6.113}
\end{align}
\begin{align}
\operatorname{supp}\nu_{(x,t)}\subseteq\phi^{-1}(\{\phi(u_r(x,t))\}). \label{align:ch6-6.114}
\end{align}
\end{enumerate}
\end{proposition}

\begin{proof}
对每个 $f\in C(\mathbb{R})$，序列 $\{(f\circ\phi)(u_{n_k})\}$ 在 $L^1(Q)$ 中有界且一致可积，故（可能抽取子列，不重新标号）由\cref{proposition:ch05-弱紧YoungL1}得
\begin{align}
f(\phi(u_{n_k}))\rightharpoonup(f\circ\phi)^* \text{ 在 }L^1(Q),\quad
(f\circ\phi)^*(x,t):=\int_\mathbb{R}f(\phi(y))\,\mathrm{d}\nu_{(x,t)}(y). \label{align:ch6-6.115}
\end{align}
注意由\eqref{align:ch6-6.87}中的收敛，有
\begin{align*}
f(\phi(u_{n_k}))\to f(\phi^*) \text{ 在 }Q\text{ 中几乎处处},
\end{align*}
从而（回忆 $\phi$ 有界）
\begin{align*}
f(\phi(u_{n_k}))\to f(\phi^*) \text{ 在 }L^1(Q).
\end{align*}
由上述收敛与\eqref{align:ch6-6.115}得
\begin{align}
f(\phi^*(x,t))=\int_\mathbb{R}f(\phi(y))\,\mathrm{d}\nu_{(x,t)}(y)
\quad\text{对所有 }(x,t)\in Q\setminus N, \label{align:ch6-6.116}
\end{align}
其中 $N\subseteq Q$ 为零测集。

固定 $(x,t)\in Q\setminus N$，置 $B=\mathbb{R}\setminus\phi^{-1}(\{\phi^*(x,t)\})$。选取 $f\in C(\mathbb{R})$，使当 $z\neq\phi^*(x,t)$ 时 $f(z)>0$，当 $z=\phi^*(x,t)$ 时 $f(z)=0$，由\eqref{align:ch6-6.116}得
\begin{align*}
0=\int_B f(\phi(y))\,\mathrm{d}\nu_{(x,t)}(y).
\end{align*}
因对所有 $y\in B$ 有 $f(\phi(y))>0$，上述等式保证 $\nu_{(x,t)}(B)=0$。因 $B=\mathbb{R}\setminus\phi^{-1}(\{\phi^*(x,t)\})$ 为开集，(i) 得证。

关于 (ii)，注意由\eqref{align:ch6-6.102a}、\eqref{align:ch6-6.103a}与\eqref{align:ch6-6.112}，
\begin{align*}
u_r(x,t)=\int_\mathbb{R}y\,\mathrm{d}\nu_{(x,t)}(y)
=\int_{\phi^{-1}(\{\phi^*(x,t)\})}y\,\mathrm{d}\nu_{(x,t)}(y)
\in\phi^{-1}(\{\phi^*(x,t)\}),
\end{align*}
因为 $\phi^{-1}(\{\phi^*(x,t)\})$ 为闭区间，$\nu_{(x,t)}$ 为概率测度。于是对几乎所有 $(x,t)\in Q$，有 $\phi(u_r(x,t))=\phi^*(x,t)$，\eqref{align:ch6-6.113}中第二个陈述随之成立，因为由\eqref{align:ch6-6.82}与\eqref{align:ch6-6.84}，$\phi^*\in L^2(0,T;H^1_0(\Omega))$。因 $\phi(u_r)=\phi^*$，由\eqref{align:ch6-6.112}得\eqref{align:ch6-6.114}。证明完成。
\end{proof}

\begin{proof}[证明\cref{theorem:ch06-6.6.1}]
设 $u\in L^\infty_{w^*}(0,T;\mathscr{R}_f(\Omega))$ 由\eqref{align:ch6-6.102b}定义。则由\cref{proposition:ch06-6.6.3}--\cref{proposition:ch06-6.6.6}，结论成立。
\end{proof}

现在可以证明\cref{theorem:ch06-6.4.1}。

\begin{proof}[证明\cref{theorem:ch06-6.4.1}]
我们证明\cref{theorem:ch06-6.6.1}所给出的测度 $u\in L^\infty_{w^*}(0,T;\mathscr{R}_f(\Omega))$ 是问题 (P) 的弱解。因 $\phi(u_r)\in L^2(0,T;H^1_0(\Omega))$，\cref{definition:ch06-弱解}的条件 (i) 满足。设 $\{u_{n_k}\}$ 为逼近问题解的子列，使 $u_{n_k}\stackrel{*}{\rightharpoonup}u$ 在 $L^\infty_{w^*}(0,T;\mathscr{R}_f(\Omega))$（见\eqref{align:ch6-6.77}）。考虑 $(Pn_k)$ 的弱形式：
\begin{align*}
\iint_Q\{u_{n_k}\partial_t\zeta-\nabla\phi_{n_k}(u_{n_k})\cdot\nabla\zeta\}\,\mathrm{d}x\,\mathrm{d}t
=-\int_\Omega u_{0n_k}(x)\zeta(x,0)\,\mathrm{d}x
\end{align*}
其中 $\zeta\in C^1([0,T];C^1_c(\Omega))$，$\zeta(\cdot,T)=0$ 在 $\Omega$ 上。令 $k\to\infty$，利用\eqref{align:ch6-6.77}、\eqref{align:ch6-6.84}与\eqref{align:ch6-6.113}、以及\eqref{align:ch6-6.49c}，得\eqref{align:ch6-6.2}。故结果第一部分成立。

以 $+$ 号为例证明 (i)、(ii) 款。为此注意由\eqref{align:ch6-6.103b}与不等式\eqref{align:ch6-6.109}，对所有 $\rho\in C([0,T];C_c(\Omega))$，$\zeta\geqslant0$ 在 $Q$ 中，$\zeta(\cdot,T)=0$ 在 $\Omega$ 上，有
\begin{align*}
-\int_0^T\langle u^+_s(\cdot,t),\zeta(\cdot,t)\rangle\,\mathrm{d}t
\leqslant\langle u^+_{0s},\zeta(\cdot,0)\rangle.
\end{align*}
在上一不等式中置 $\zeta(x,t)=\rho(x)h(t)$，$\rho\in C_c(\Omega)$，$\rho\geqslant0$。如\cref{theorem:ch06-6.4.5}证明中那样，通过适当选取函数 $h$，(i)、(ii) 两款均容易得到。
\end{proof}

设 $\pm\phi_{\pm\infty}>0$，设 $f_{c,\pm}:[0,\pm\phi_{\pm\infty}]\to[0,1]$ 及 $f_c\in C^1([0,\pm\phi_{\pm\infty}])$ 满足
\begin{align}
f_{c,\pm}(y)=
\begin{cases}
0 & \text{若 }y\in[0,c],\\
1 & \text{若 }y\in[\frac{c\pm\phi_{\pm\infty}}{2},\pm\phi_{\pm\infty}]
\end{cases} \label{align:ch6-6.117}
\end{align}
对某个 $c\in(0,\pm\phi_{+\infty})$。为证\cref{theorem:ch06-6.4.2}，需要如下引理。

\begin{lemma}\label{lemma:ch06-6.6.7}
设 (A0)--(A2) 成立，$\pm\phi_{\pm\infty}>0$，$u$ 为\cref{theorem:ch06-6.4.1}所给出的问题 (P) 弱解。则对几乎所有 $t\in(0,T)$，对任意 $c\in(0,\pm\phi_{\pm\infty})$，及任意 $\rho\in C^1_c(\Omega)$，$\rho\geqslant0$，有
\begin{align}
\langle[u^\pm_s(\cdot,t)]_{d,2},\rho\rangle
\leqslant\langle[u^\pm_{0s}]_{d,2},f_c(\phi^\pm(u_r(\cdot,t)))\rho\rangle. \label{align:ch6-6.118}
\end{align}
\end{lemma}

\begin{proof}[证明\cref{theorem:ch06-6.4.2}]
我们仅证\eqref{align:ch6-6.21a}中 $+$ 号等式，$-$ 号证明类似。注意 $\phi_{+\infty}=0$ 蕴含对所有 $y\in[0,\infty)$ 有 $\phi(y)=0$，因而在此情形 $\phi^+\equiv0$，$\mathscr{S}^+_t=\Omega$，\eqref{align:ch6-6.2}显然满足。

设 $\phi_{+\infty}>0$，为简便置 $f_c\equiv f_{c,+}$（见\eqref{align:ch6-6.117}）。固定任意 $t\in(0,T)$，使对所有非负 $\rho\in C^1_c(\Omega)$，
\begin{align*}
\langle[u^+_s(\cdot,t)]_{d,2},\rho\rangle
\leqslant\langle[u^+_{0s}]_{d,2},f_c(\phi^+(u_r(\cdot,t)))\rho\rangle,
\end{align*}
设 $g_c\in C^1(\mathbb{R})$ 为任意函数，$g_c\geqslant0$，$\operatorname{supp}g_c\subseteq[0,c)$。由 $f_c$ 的定义（见\eqref{align:ch6-6.117}），在 $[0,\infty)$ 上 $f_c g_c=0$。

回忆对几乎所有 $t\in(0,T)$，函数 $\phi^+(u_r(\cdot,t))$ 等同于其 $2$-拟连续代表元，后者定义到相差 $2$-零容集（见\cref{remark:ch06-6.3.1}）。故对几乎所有 $t\in(0,T)$ 及任意 $c\in(0,\phi_{+\infty})$，有
\begin{align}
f_c(\phi^+(u_r(x,t)))g_c(\phi^+(u_r(x,t)))=0
\quad\text{对所有 }x\in\Omega^*, \label{align:ch6-6.119}
\end{align}
对某个子集 $\Omega^*\subseteq\Omega$，$C_2(\Omega\setminus\Omega^*)=0$。

因 $g_c((\phi^+(u_r))(\cdot,t))\in H^1(\Omega)\cap L^\infty(\Omega)$，函数 $g_c((\phi^+(u_r))(\cdot,t))\widetilde\rho$（$\widetilde\rho\in C^1_c(\Omega)$，$\widetilde\rho\geqslant0$）属于 $H^1_0(\Omega)\cap L^\infty(\Omega)$，故可作为不等式\eqref{align:ch6-6.118}中的试验函数。由\eqref{align:ch6-6.119}得
\begin{align}
\begin{aligned}
&\langle[u_s(x,t)]_{d,2},g_c(\phi^+(u_r(x,t)))\widetilde\rho\rangle\\
&\quad\leqslant\langle[u^\pm_{0s}]_{d,2},f_c(\phi^\pm(u_r(x,t)))g_c(\phi^+(u_r(x,t)))\widetilde\rho\rangle=0
\end{aligned} \label{align:ch6-6.120}
\end{align}
对每个如上 $g_c$ 与 $\widetilde\rho$ 及 $x\in\Omega^*$ 成立。

取 $c=c_q:=\phi_{+\infty}-\frac{1}{q}$，使 $\operatorname{supp}g_{c_q}\subseteq[0,\phi_{+\infty}-\frac{1}{q})$（$q\in\mathbb{N}$）。显然这蕴含
\begin{align}
[u_s(\cdot,t)]_{d,2}(E^t_q)=0, \label{align:ch6-6.121}
\end{align}
其中
\begin{align*}
E^t_q:=\{x\in\Omega^*\mid\phi^+(u_r(x,t))\leqslant\phi_{+\infty}-\frac{2}{q}\}.
\end{align*}
因对所有 $q\in\mathbb{N}$ 有 $E^t_q\subseteq E^t_{q+1}$，且
\begin{align*}
F^t:=\{x\in\Omega^*\mid\phi^+(u_r(x,t))<\phi_{+\infty}\}=\bigcup_{q=1}^\infty E^t_q,
\end{align*}
由等式\eqref{align:ch6-6.121}得
\begin{align*}
[u_s(\cdot,t)]_{d,2}(F^t)=\lim_{q\to\infty}[u_s(\cdot,t)]_{d,2}(E^t_q)=0.
\end{align*}
于是
\begin{align}
[u_s(\cdot,t)]_{d,2}\llcorner\Omega^*
=[u_s(\cdot,t)]_{d,2}\llcorner\{x\in\Omega^*\mid\phi^+(u_r(x,t))=\phi_{+\infty}\}. \label{align:ch6-6.122}
\end{align}
另一方面，$[u_s(\cdot,t)]_{d,2}(\Omega\setminus\Omega^*)=0$，因 $C_2(\Omega\setminus\Omega^*)=0$。由\eqref{align:ch6-6.122}，结论成立。
\end{proof}

\begin{proof}[证明\cref{lemma:ch06-6.6.7}]
我们仅证\eqref{align:ch6-6.21a}中 $+$ 号情形，另一情形证明类似。设 $\phi_{+\infty}\in(0,\infty)$。对每个 $m\in\mathbb{N}$，置
\begin{align*}
h_m(y):=(y-m)\chi_{[m,m+1]}(y)+\chi_{(m+1,+\infty)}(y),\quad
H_m(y):=\int_0^y h_m(z)\,\mathrm{d}z \quad(y\in\mathbb{R}).
\end{align*}
设 $\{u_{n_k}\}$ 与 $u$ 由\cref{theorem:ch06-6.6.1}给出。因 $\|h_m\|_\infty=1$，由\eqref{align:ch6-6.54}与\eqref{align:ch6-6.64}，对所有 $\rho\in C^1_c(\Omega)$ 及 $m\in\mathbb{N}$，序列 $\{\int_\Omega H_m(u_{n_k}(x,\cdot))\rho(x)\,\mathrm{d}x\}$ 在 $BV(0,T)$ 中有界。于是存在零测集 $N_1\subseteq(0,T)$ 及 $\{u_{n_k}\}$ 的子列（为简便不重新标号），使序列 $\{\int_\Omega H_m(u_{n_k}(x,t))\rho(x)\,\mathrm{d}x\}$ 对所有 $t\in(0,T)\setminus N_1$ 在 $L^1((0,T))$ 中收敛。由 $C_0(\Omega)$ 的可分性及标准对角论证，零测集 $N_1$ 可取得与 $\rho\in C^1_c(\Omega)$ 及 $m\in\mathbb{N}$ 均无关。因 $\{\int_\Omega H_m(u_{n_k}(x,t))\rho(x)\,\mathrm{d}x\}$ 对 $t\in(0,T)\setminus N_1$ 收敛，因 $\lim_{y\to+\infty}H_m(y)/y=1$，$\lim_{y\to-\infty}H_m(y)/y=0$，在\eqref{align:ch6-6.98}中取 $\zeta(x,t)=\rho(x)h(t)$ 并利用\eqref{align:ch6-6.103b}，对所有 $t\in(0,T)\setminus N_1$，$\rho\in C^1_c(\Omega)$，及 $m\in\mathbb{N}$，得
\begin{align}
\lim_{k\to\infty}\int_\Omega H_m(u_{n_k}(x,t))\rho\,\mathrm{d}x
=\int_\Omega\left(\int_\mathbb{R}H_m(y)\,\mathrm{d}\nu_{(x,t)}(y)\right)\rho(x)\,\mathrm{d}x
+\langle u^+_s(\cdot,t),\rho\rangle. \label{align:ch6-6.123}
\end{align}

由\eqref{align:ch6-6.88}--\eqref{align:ch6-6.89}与\eqref{align:ch6-6.113}，存在零测集 $N_2\subseteq(0,T)$，使对所有 $t\in(0,T)\setminus N_2$，
\begin{align}
\phi(u_{n_k}(\cdot,t))\to\phi^+(u_r(\cdot,t)) \quad\text{在 }\Omega\text{ 中几乎处处}, \label{align:ch6-6.124a}
\end{align}
\begin{align}
\phi(u_{n_k}(\cdot,t))\rightharpoonup\phi^+(u_r(\cdot,t)) \quad\text{在 }H^1_0(\Omega). \label{align:ch6-6.124b}
\end{align}
置 $N:=N_1\cup N_2$，固定 $t_0\in(0,T)\setminus N$。注意对每个充分大的 $m\in\mathbb{N}$，有
\begin{align*}
H_m(s)[1-f_c(\phi(s))]=0
\end{align*}
对每个 $s\in\mathbb{R}$ 成立，其中 $f_c\equiv f_{c,+}$ 由\eqref{align:ch6-6.117}给出。于是对所有充分大的 $m\in\mathbb{N}$ 及 $\rho\in C^1_c(\Omega)$，$\rho\geqslant0$，由\eqref{align:ch6-6.60}、\eqref{align:ch6-6.62}、\eqref{align:ch6-6.74}及\eqref{align:ch6-6.124a}，有
\begin{align*}
\begin{aligned}
&\int_\Omega H_m(u_{n_k}(x,t_0))\rho(x)\,\mathrm{d}x\\
&\quad=\int_\Omega H_m(u_{n_k}(x,t_0))f_c(\phi(u_{n_k}(x,t_0)))\rho(x)\,\mathrm{d}x\\
&\quad\leqslant\left|\iint_{\{u_{n_k}>m\}}h_m(u_{n_k})\nabla\phi_{n_k}(u_{n_k})\cdot
\nabla(f_c(\phi(u_{n_k}(x,t_0)))\rho)\,\mathrm{d}x\,\mathrm{d}t\right|\\
&\qquad+\int_\Omega H_m(u_{0n_k})f_c(\phi(u_{n_k}(x,t_0)))\rho(x)\,\mathrm{d}x\\
&\quad\leqslant\|f_c(\phi(u_{n_k}(\cdot,t_0)))\rho\|_{H^1_0(\Omega)}
\left(\iint_{\{u_{n_k}>m\}}|\nabla\phi_{n_k}(u_{n_k})|^2\,\mathrm{d}x\,\mathrm{d}t\right)^{1/2}\\
&\qquad+\int_\Omega H_m(u_{0n_k})f_c(\phi(u_{n_k}(x,t_0)))\rho(x)\,\mathrm{d}x\\
&\quad\leqslant M_c[\phi_{+\infty}-\phi(m)+n_k^{-1/2}]^{1/2}\\
&\qquad+\int_\Omega[u^+_0]_{n_k}f_c(\phi(u_{n_k}(x,t_0)))\rho(x)\,\mathrm{d}x\\
&\quad\leqslant M_c[\phi_{+\infty}-\phi(m)+n_k^{-1/2}]^{1/2}\\
&\qquad+\int_\Omega[\delta^+]_{n_k}f_c(\phi(u_{n_k}(x,t_0)))\rho(x)\,\mathrm{d}x
+\int_\Omega[\gamma^+]_{n_k}\rho(x)\,\mathrm{d}x.
\end{aligned}
\end{align*}
上述常数 $M_c$ 取得使 $M_0\|f_c(\phi(u_{n_k}(\cdot,t_0)))\rho\|_{H^1_0(\Omega)}\leqslant M_c$ 对所有 $k\in\mathbb{N}$，其中 $M_0$ 为\eqref{align:ch6-6.75}中的常数。由\eqref{align:ch6-6.49d}中第二个收敛、\eqref{align:ch6-6.49e}及\eqref{align:ch6-6.123}--\eqref{align:ch6-6.124a}，先令 $k\to+\infty$ 再令 $m\to+\infty$，对所有 $c\in(0,\phi_{+\infty})$ 及如上 $\rho$，得
\begin{align}
\langle u^+_s(\cdot,t_0),\rho\rangle
\leqslant\langle[u^+_0]_{d,2},f_c(\phi^+(u_r(\cdot,t_0)))\rho\rangle
+\langle[u^+_0]_{c,2},\rho\rangle. \label{align:ch6-6.125}
\end{align}
由\cref{theorem:ch06-6.2.1}（见\eqref{align:ch6-6.12}）与等式\eqref{align:ch03-3.82}，对几乎所有 $t\in(0,T)$，有 $[u^\pm_s(\cdot,t)]_{c,2}=[u^\pm(\cdot,t)]_{c,2}=[u^\pm_0]_{c,2}$。于是不等式\eqref{align:ch6-6.125}给出\eqref{align:ch6-6.118}。结果得证。
\end{proof}

\begin{proof}[证明\cref{proposition:ch06-6.4.4}]
由\cref{definition:ch06-构造解}，(P) 的每个构造解 $u$ 是逼近问题解序列子列 $\{u_{n_k}\}$ 在 $L^\infty_{w^*}(0,T;\mathscr{M}(\Omega))$ 中的弱$^*$极限（见\cref{theorem:ch06-6.4.1}）。在\eqref{align:ch6-6.61}中取 $t_1=0$，$t_2=T$，得
\begin{align}
\iint_Q\{F_\phi(u_{n_k})\partial_t\zeta-f(\phi(u_{n_k}))\nabla\phi_{n_k}(u_{n_k})\cdot\nabla\zeta\}\,\mathrm{d}x\,\mathrm{d}t
\geqslant-\int_\Omega F_\phi(u_{0n_k}(x))\zeta(x,0)\,\mathrm{d}x \label{align:ch6-6.126}
\end{align}
对所有 $\zeta\in C^1([0,T];C^1_c(\Omega))$，$\zeta\geqslant0$ 在 $Q$ 中，$\zeta(\cdot,T)=0$ 在 $\Omega$ 上，及所有非减 $f\in C(\mathbb{R})$ 成立，函数 $F_\phi$ 由\eqref{align:ch6-6.9}定义。

由\eqref{align:ch6-6.84}、\eqref{align:ch6-6.87}中的收敛及\eqref{align:ch6-6.113}中的等式，直接得
\begin{align}
\lim_{k\to\infty}\iint_Q f(\phi(u_{n_k}))\nabla\phi_{n_k}(u_{n_k})\cdot\nabla\zeta\,\mathrm{d}x\,\mathrm{d}t
=\iint_Q f(\phi(u_r))\nabla\phi(u_r)\cdot\nabla\zeta\,\mathrm{d}x\,\mathrm{d}t, \label{align:ch6-6.127}
\end{align}
而由\eqref{align:ch6-6.51}，
\begin{align}
\begin{aligned}
&\lim_{k\to\infty}\int_\Omega F_\phi(u_{0n_k}(x))\zeta(x,0)\,\mathrm{d}x\\
&\quad=\int_\Omega F_\phi(u_{0r})\zeta(x,0)\,\mathrm{d}x
+M^+_{F_\phi}\langle u^+_{0s},\zeta(\cdot,0)\rangle
-M^-_{F_\phi}\langle u^-_{0s},\zeta(\cdot,0)\rangle
\end{aligned} \label{align:ch6-6.128}
\end{align}
其中 $M^\pm_{F_\phi}\in\mathbb{R}$ 如\eqref{align:ch6-6.10}。另一方面，由\cref{proposition:ch06-6.6.4}与\cref{proposition:ch06-6.6.5}，有
\begin{align}
\begin{aligned}
&\lim_{k\to\infty}\iint_Q F_\phi(u_{n_k})\partial_t\zeta\,\mathrm{d}x\,\mathrm{d}t\\
&\quad=\iint_Q F^*_\phi\partial_t\zeta\,\mathrm{d}x\,\mathrm{d}t
+M^+_{F_\phi}\int_0^T\langle u^+_s(\cdot,t),\partial_t\zeta(\cdot,t)\rangle_\Omega\,\mathrm{d}t\\
&\qquad-M^-_{F_\phi}\int_0^T\langle u^-_s(\cdot,t),\partial_t\zeta(\cdot,t)\rangle_\Omega\,\mathrm{d}t
\end{aligned} \label{align:ch6-6.129}
\end{align}
（见\eqref{align:ch6-6.102a}），其中对几乎所有 $(x,t)\in Q$，
\begin{align*}
F^*_\phi(x,t)=\int_\mathbb{R}F_\phi(y)\,\mathrm{d}\nu_{(x,t)}(y),
\end{align*}
Young 测度 $\nu\in\mathcal{Y}^+(Q;\mathbb{R})$ 由\cref{proposition:ch06-6.6.4}给出。由\eqref{align:ch6-6.114}有
\begin{align}
\begin{aligned}
\int_\mathbb{R}F_\phi(y)\,\mathrm{d}\nu_{(x,t)}(y)
&=\int_{\operatorname{supp}\nu_{(x,t)}}F_\phi(y)\,\mathrm{d}\nu_{(x,t)}(y)\\
&=\int_{\operatorname{supp}\nu_{(x,t)}}\mathrm{d}\nu_{(x,t)}(y)\int_0^y f(\phi(z))\,\mathrm{d}z
=F_\phi(u_r(x,t))
\end{aligned} \label{align:ch6-6.130}
\end{align}
对几乎所有 $(x,t)\in Q$ 成立。

由上述说明，在\eqref{align:ch6-6.126}中令 $k\to\infty$，利用\eqref{align:ch6-6.127}--\eqref{align:ch6-6.130}，证明 $u$ 满足熵不等式\eqref{align:ch6-6.11}。结论成立。
\end{proof}

\section{$\phi$ 无界的情形}\label{section:ch06-无界phi}

本节致力于在更一般假设下研究问题 (P)，特别允许函数 $\phi$ 无界。除\cref{theorem:ch6-正则化}的简短证明外，我们仅描述主要结果并给出证明梗概（完整证明请感兴趣的读者参见文献~\cite{ref87}）。

此后假设 (A1)--(A2) 由如下两条假设替代：
\begin{align*}
\phi\in C(\mathbb{R}) \text{ 在 }\mathbb{R}\text{ 上非减且非常数，}\phi(0)=0; \qquad(A1')
\end{align*}
\begin{align}
\text{存在 }M_0>0\text{ 与 }\alpha\in[1,\infty)\text{ 使 }|\phi(y)|\leqslant M_0(1+|y|^\alpha)
\quad(y\in\mathbb{R}). \qquad(A2') \label{align:ch6-A2p}
\end{align}

由假设 (A1')，存在
\begin{align}
\lim_{y\to\pm\infty}\phi(y)=:\phi_{\pm\infty}\in\mathbb{R}. \label{align:ch6-6.131}
\end{align}
$\phi_{\pm\infty}$ 有限或无穷两种情形均允许，且要么 $\phi_{-\infty}\leqslant0<\phi_{+\infty}$，要么 $\phi_{-\infty}<0\leqslant\phi_{+\infty}$。对 $\phi$ 的唯一限制是其增长至多为幂次；不使用下界估计（除\cref{proposition:ch06-6.7.9}中的正则化结果外），不对扩散率 $\phi'$ 的存在性作假设。

\subsection{解的定义}\label{subsection:ch06-无界解定义}

记 $T_K(y):=\max\{-K,\min\{y,K\}\}$（$y\in\mathbb{R}$，$K>0$）为标准截断函数。

\begin{definition}\label{definition:ch06-无界弱解}
设 (A0) 成立。问题 (P) 的弱解指任意 $u\in L^\infty_{w^*}(0,T;\mathscr{R}_f(\Omega))$，满足：
\begin{enumerate}[(i)]
\item $\phi(u_r)\in L^q(0,T;W^{1,q}_0(\Omega))$ 对任意 $q\in[1,1+\frac{1}{\alpha N+1})$，且 $T_K(\phi(u_r))\in L^2(0,T;H^1_0(\Omega))$ 对任意 $K>0$；
\item 对每个 $\zeta\in C^1([0,T];C^1_c(\Omega))$，$\zeta(\cdot,T)=0$ 在 $\Omega$ 上，有
\begin{align}
\int_0^T\langle u(\cdot,t),\partial_t\zeta(\cdot,t)\rangle\,\mathrm{d}t
= \iint_Q\nabla\phi(u_r)\cdot\nabla\zeta\,\mathrm{d}x\,\mathrm{d}t
-\langle u_0,\zeta(\cdot,0)\rangle. \label{align:ch6-6.132}
\end{align}
\end{enumerate}
\end{definition}

易见在此情形等式\eqref{align:ch6-6.5a}同样成立，即初值条件在 $\mathscr{R}_f(\Omega)$ 的弱$^*$意义下满足。

\begin{remark}\label{remark:ch06-6.7.1}
(i) 因 $[T_K(\phi)]^\pm=T_K(\phi^\pm)$，对 (P) 的每个弱解，若 $\pm\phi_{\pm\infty}<+\infty$，则 $\phi^\pm(u_r)\in L^2(0,T;H^1_0(\Omega))$。故在此情形，\cref{definition:ch06-无界弱解}化归为\cref{definition:ch06-弱解}。

(ii) 若 $\phi$ 有界，则 (P) 的每个弱解 $u$ 具有如下性质：
\begin{align}
\phi(u_r)\in L^2(0,T;H^1_0(\Omega)), \label{align:ch6-6.133a}
\end{align}
\begin{align}
u(\cdot,t)-u_0\in H^{-1}(\Omega) \quad\text{对几乎所有 }t\in(0,T), \label{align:ch6-6.133b}
\end{align}
\begin{align}
\operatorname*{ess\,lim}_{t\to0^+}\|u(\cdot,t)-u_0\|_{H^{-1}(\Omega)}=0. \label{align:ch6-6.133c}
\end{align}
事实上，\eqref{align:ch6-6.133a}已在 (i) 中指出，而由等式\eqref{align:ch6-6.132}得
\begin{align}
u(\cdot,t)-u_0=\Delta\left(\int_0^t\phi(u_r(\cdot,s))\,\mathrm{d}s\right)
\quad\text{在 }\mathscr{D}'(\Omega)\text{ 中对几乎所有 }t\in(0,T). \label{align:ch6-6.134}
\end{align}
由\eqref{align:ch6-6.133a}，等式\eqref{align:ch6-6.134}给出
\begin{align*}
\|u(\cdot,t)-u_0\|^2_{H^{-1}(\Omega)}
=\int_\Omega\left|\nabla\int_0^t\phi(u_r(x,s))\,\mathrm{d}s\right|^2\mathrm{d}x
\leqslant\int_0^t\int_\Omega|\nabla\phi(u_r)|^2\,\mathrm{d}x\,\mathrm{d}s,
\end{align*}
从而\eqref{align:ch6-6.133b}与\eqref{align:ch6-6.133c}成立。断言得证。
\end{remark}

\subsection{持续性与唯一性}\label{subsection:ch06-无界持续唯一}

$\phi$ 在无穷远次线性时的持续性性质（与\cref{theorem:ch06-6.2.1}类似）由如下结果给出。

\begin{theorem}\label{theorem:ch6-正则化}
设 (A0) 与 (A1') 成立，且存在 $M>0$ 与 $\alpha\in(0,\frac{N-2}{N}]$，使
\begin{align*}
|\phi(y)|\leqslant M(1+|y|)^\alpha \quad(y\in\mathbb{R}).
\end{align*}
设 $u$ 为问题 (P) 的弱解。则存在零测集 $F\subseteq(0,T)$，使对每个 $t\in(0,T)\setminus F$，$u(\cdot,t)$ 的 $C_{2,\frac{1}{1-\alpha}}$-集中部分在时间上为常数：
\begin{align}
[u(\cdot,t)]_{c,2,\frac{1}{1-\alpha}}=[u_0]_{c,2,\frac{1}{1-\alpha}}. \label{align:ch6-6.135}
\end{align}
\end{theorem}

\begin{remark}\label{remark:ch06-6.7.2}
注意\cref{theorem:ch6-正则化}与文献~\cite{ref26,ref80}中关于快扩散多孔介质方程 Cauchy 问题的结果一致。特别地，设 $u_0\in\mathscr{R}^+_{c,2}(\mathbb{R}^N)\cap\mathscr{R}^+_{d,2,\frac{1}{1-\alpha}}(\mathbb{R}^N)$（回忆由\cref{proposition:ch03-容量包含}，$\mathscr{R}_{c,2,\frac{1}{1-\alpha}}(\mathbb{R}^N)\subseteq\mathscr{R}_{c,2}(\mathbb{R}^N)$）。则文献~\cite{ref26,ref80}意义下的解满足对所有 $t\in(0,T)$，$u(\cdot,t)\in L^1(\mathbb{R}^N)$，故在此情形等式\eqref{align:ch6-6.12}不成立。这提示了\cref{theorem:ch6-正则化}的如下解释：当无穷远附近扩散率更强时，只有关于更强容量集中的测度在时间上保持常数。

在此联系中注意，对任意 $\alpha,\beta\in(0,\frac{N-2}{N}]$，$\alpha<\beta$：
\begin{enumerate}[(a)]
\item $C_{2,\frac{1}{1-\beta}}$-容量强于 $C_{2,\frac{1}{1-\alpha}}$-容量；
\item 关于 $C_{2,\frac{1}{1-\beta}}$-容量集中的测度也关于 $C_{2,\frac{1}{1-\alpha}}$-容量集中（见\cref{theorem:ch03-容量单调}与\cref{proposition:ch03-容量包含}）。
\end{enumerate}
进而，关于 $C_{2,\frac{1}{1-\alpha}}$-容量（$\alpha\in(0,\frac{N-2}{N}]$）集中的测度也关于 Newton 容量集中（见\cref{proposition:ch03-容量包含}）。

注意对 $\alpha\in(\frac{N-2}{N},1)$，不存在\cref{theorem:ch6-正则化}的类似结果，因为在此情形每个测度关于 $C_{2,\frac{1}{1-\alpha}}$-容量都是弥散的（见\cref{proposition:ch03-点容量}）。
\end{remark}

\begin{proof}[证明\cref{theorem:ch6-正则化}]
设 $u$ 为问题 (P) 的弱解。由\eqref{align:ch6-6.132}，对每个 $\rho\in C^2_c(\Omega)$ 及 $h\in C^1([0,T])$，$h(0)=1$，$h(T)=0$，得
\begin{align*}
\int_0^T\langle u(\cdot,t),\rho\rangle h'(t)\,\mathrm{d}t
=-\iint_Q\phi(u_r)\Delta\rho\,h(t)\,\mathrm{d}x\,\mathrm{d}t-\langle u_0,\rho\rangle.
\end{align*}
由标准逼近论证，由上述等式对几乎所有 $t\in(0,T)$ 得
\begin{align}
|\langle u(\cdot,t)-u_0,\rho\rangle|
=\left|\int_0^t\int_\Omega\phi(u_r)\Delta\rho\,\mathrm{d}x\,\mathrm{d}t\right|
\leqslant Mt^{1-\alpha}\left(\iint_Q(1+|u_r|)\,\mathrm{d}x\,\mathrm{d}t\right)^\alpha
\left(\int_\Omega|\Delta\rho|^{\frac{1}{1-\alpha}}\,\mathrm{d}x\right)^{1-\alpha}. \label{align:ch6-6.136}
\end{align}
由上述不等式，存在零测集 $F\subseteq(0,T)$，使对所有 $t\in(0,T)\setminus F$，$u(\cdot,t)-u_0\in W^{-2,\frac{1}{\alpha}}(\Omega)=(W^{2,\frac{1}{1-\alpha}}_0(\Omega))^*$。由文献~\cite{ref10} 引理 4.1，这蕴含对任意此类 $t$，测度 $u(\cdot,t)-u_0$ 是 $C_{2,\frac{1}{1-\alpha}}$-弥散的，因而 $u(\cdot,t)-u_0$ 的 $C_{2,\frac{1}{1-\alpha}}$-集中部分为零。故等式\eqref{align:ch6-6.135}成立。
\end{proof}

如前所述，解的持续性性质与唯一性相联系，需要附加条件来确定问题 (P) 的唯一类。在此情形，这些条件由如下定义给出，它推广了\cref{definition:ch06-相容性条件}。

\begin{definition}\label{definition:ch06-无界相容性}
设 (A0) 与 (A1') 成立。问题 (P) 的弱解 $u$ 满足相容性条件，若对几乎所有 $t\in(0,T)$：
\begin{enumerate}[(i)]
\item 若 $\phi_{\pm\infty}\in\mathbb{R}$，则条件\eqref{align:ch6-6.21a}满足；
\item 若 $\phi_{\pm\infty}=\pm\infty$，则
\begin{align}
[u^\pm_s(\cdot,t)]_{d,2}=0. \label{align:ch6-6.137}
\end{align}
\end{enumerate}
若 $\phi_{\pm\infty}\in\mathbb{R}$，\cref{remark:ch06-6.3.1}在此情形同样适用于 $\phi^\pm(u_r)$。
\end{definition}

相容性条件的重要性由如下唯一性结果凸显，由\cref{remark:ch06-6.7.1}(ii)，它将\cref{theorem:ch06-6.3.1}作为特例包含。

\begin{theorem}\label{theorem:ch06-6.7.2}
设 (A0) 与 (A1') 成立。设 $u_1$ 与 $u_2$ 为问题 (P) 的两个弱解，满足相容性条件，且
\begin{align}
\phi(u_{1r})-\phi(u_{2r})\in L^2(\tau,T;H^1_0(\Omega)) \quad\text{对所有 }\tau\in(0,T), \label{align:ch6-6.138a}
\end{align}
\begin{align}
u_1(\cdot,t)-u_2(\cdot,t)\in H^{-1}(\Omega) \quad\text{对几乎所有 }t\in(0,T), \label{align:ch6-6.138b}
\end{align}
\begin{align}
\operatorname*{ess\,lim}_{t\to0^+}\|u_1(\cdot,t)-u_2(\cdot,t)\|_{H^{-1}(\Omega)}=0. \label{align:ch6-6.138c}
\end{align}
则 $u_1=u_2$ 在 $L^\infty_{w^*}(0,T;\mathscr{R}_f(\Omega))$。
\end{theorem}

\subsection{存在性}\label{subsection:ch06-无界存在}

同样地，我们遵循\cref{section:ch06-存在正则}中概述的构造性存在方法，故第一个存在性结果是\cref{theorem:ch06-6.4.1}的对应物。

\begin{theorem}\label{theorem:ch06-6.7.3}
设 (A0) 与 (A1')--(A2') 成立。则存在问题 (P) 的弱解，它作为逼近问题 (Pn) 解序列在 $\mathscr{R}_f(Q)$ 的弱$^*$拓扑下的极限点而得到。
\end{theorem}

为证明满足相容性条件的弱解存在性，现在使用两步过程，它对 $u_0$ 采用如下附加假设：
\begin{align*}
[u^\pm_0]_{d,2}\in H^{-1}(\Omega), \qquad(A3)
\end{align*}
\begin{align*}
[u^\pm_0]_{c,2}=[u^\pm_0]_{c,2}\llcorner K^\pm_0,\quad
K^\pm_0\subseteq\Omega\text{ 紧且 }C_2(K^\pm_0)=0. \qquad(A4)
\end{align*}

首先证明如下结果（见文献~\cite{ref87} 命题 3.4）。

\begin{proposition}\label{proposition:ch06-6.7.4}
设 (A0) 与 (A1')--(A2') 成立。若 $\phi_{\pm\infty}=\pm\infty$，则设 $u^\pm_0$ 还满足 (A3)--(A4)。则存在问题 (P) 满足相容性条件的弱解。
\end{proposition}

由\cref{proposition:ch06-6.7.4}，若 (A1')--(A2') 成立且 $\phi$ 有界，则存在 (P) 满足\eqref{align:ch6-6.21a}的弱解，这是已由\cref{proposition:ch06-6.4.2}所知的结果。若 $\phi$ 无界，则需进一步逼近过程以去掉假设 (A3)--(A4)。对每个 $i,j\in\mathbb{N}$，考虑问题
\begin{align*}
\begin{cases}
\partial_t u_{i,j}=\Delta\phi(u_{i,j}) & \text{在 }Q,\\
u_{i,j}=0 & \text{在 }\Gamma,\\
u_{i,j}=u_{0,i,j} & \text{在 }\Omega\times\{0\},
\end{cases}
\qquad(Pi,j)
\end{align*}
其中初值 $u_{0,i,j}\in\mathscr{R}_f(\Omega)$ 在 $\phi_{\pm\infty}=\pm\infty$ 时满足条件 (A3)--(A4)，关于 $i,j$ 具有合适的单调性，并当 $i,j\to\infty$ 时在合适意义下收敛于 $u_0$（见文献~\cite{ref87} 第 7 节；此类序列的存在性由文献~\cite{ref10} 引理 4.2 保证）。由\cref{proposition:ch06-6.7.4}，对所有 $i,j\in\mathbb{N}$，存在 (Pi,j) 满足相容性条件的弱解 $u_{i,j}$。由单调性方法证明当 $i,j\to\infty$ 时 $u_{i,j}$ 在 $L^\infty_{w^*}(0,T;\mathscr{R}_f(\Omega))$ 中的极限存在且为 (P) 满足相容性条件的弱解。于是有如下结果。

\begin{theorem}\label{theorem:ch06-6.7.5}
设 (A0) 与 (A1')--(A2') 成立。则存在问题 (P) 满足相容性条件的弱解。
\end{theorem}

现在称\cref{theorem:ch06-6.7.5}所给出的问题 (P) 弱解为构造解。若 $\phi$ 有界，则构造解由\cref{proposition:ch06-6.7.4}给出，此定义与\cref{definition:ch06-构造解}一致。我们得出结论，在此情形构造解也满足相容性条件。

由\cref{theorem:ch06-6.7.2}与\cref{theorem:ch06-6.7.5}，有如下结果，它推广了\cref{theorem:ch06-6.4.3}（回忆由\cref{remark:ch03-H负1测度}，测度 $\mu\in\mathscr{R}_f(\Omega)$ 可属于对偶空间 $H^{-1}(\Omega)=(H^1_0)^*$）。

\begin{theorem}\label{theorem:ch06-6.7.6}
设 (A0) 与 (A1')--(A2') 成立。
\begin{enumerate}[(i)]
\item 设 $\phi$ 有界。则对每个 $u_0\in\mathscr{R}_f(\Omega)$，存在问题 (P) 满足\eqref{align:ch6-6.21a}的唯一弱解 $u$。
\item 设 $\phi$ 无界。则对每个 $u_0\in\mathscr{R}_f(\Omega)\cap H^{-1}(\Omega)$，存在问题 (P) 满足相容性条件的唯一弱解 $u$。进而，
\begin{align}
\phi(u_r)\in L^2(\tau,T;H^1_0(\Omega)) \quad\text{对所有 }\tau\in(0,T), \label{align:ch6-6.139a}
\end{align}
\begin{align}
u(\cdot,t)\in H^{-1}(\Omega) \quad\text{对几乎所有 }t\in(0,T), \label{align:ch6-6.139b}
\end{align}
\begin{align}
\operatorname*{ess\,lim}_{t\to0^+}\|u(\cdot,t)-u_0\|_{H^{-1}(\Omega)}=0. \label{align:ch6-6.139c}
\end{align}
\end{enumerate}
\end{theorem}

\subsection{正则化}\label{subsection:ch06-无界正则}

同样，在此情形构造解的奇异部分既不能自发出现，也不能随时间增长。

\begin{theorem}\label{theorem:ch06-6.7.7}
设 (A0) 与 (A1')--(A2') 成立，$u$ 为问题 (P) 的构造解。则对几乎所有 $t_1,t_2\in(0,T)$，$t_1<t_2$，不等式\eqref{align:ch6-6.33}成立。
\end{theorem}

最后提及两个 $\mathscr{R}_f$--$L^1$ 正则化效应（见文献~\cite{ref87} 注 3.7 与命题 3.11）。

\begin{proposition}\label{proposition:ch06-6.7.8}
设 (A0) 与 (A1')--(A2') 成立且 $\phi_{\pm\infty}=\pm\infty$。则对每个具有 $2$-弥散初值的问题 (P) 构造解，发生瞬时正则化。
\end{proposition}

\begin{proposition}\label{proposition:ch06-6.7.9}
设 (A0) 与 (A1')--(A2') 成立，$u$ 为问题 (P) 的构造解。设存在 $\alpha_\pm>\frac{N-2}{N}$ 及 $a_\pm,b_\pm>0$，使
\begin{align}
\pm\phi(z)\geqslant a_\pm|z|^{\alpha_\pm}-b_\pm \quad\text{当 }\pm z\geqslant0. \label{align:ch6-6.140}
\end{align}
则 $u^\pm\in L^\infty(0,T;L^1(\Omega))$。
\end{proposition}

\end{document}
